\documentclass[11pt,a4paper]{article}
\usepackage{fontspec,amsmath,amssymb,amsthm,microtype}
\setmainfont{lmroman10-regular.otf}[BoldFont=lmroman10-bold.otf,ItalicFont=lmroman10-italic.otf,BoldItalicFont=lmroman10-bolditalic.otf]
\usepackage{xcolor}
\usepackage[most]{tcolorbox}
\usepackage{mathrsfs}
\providecommand{\lt}{<}
\providecommand{\gt}{>}
\setmonofont{DejaVuSansMono.ttf}[BoldFont=DejaVuSansMono-Bold.ttf,ItalicFont=DejaVuSansMono-Oblique.ttf,BoldItalicFont=DejaVuSansMono-BoldOblique.ttf,Scale=0.86]
\newfontfamily\leanfallback{STIXTwoMath-Regular.otf}[Scale=0.86]
\definecolor{leanrule}{gray}{0.62}
\definecolor{leanback}{gray}{0.962}
\definecolor{leantitle}{gray}{0.22}
\definecolor{leanmuted}{gray}{0.38}
\newtcolorbox{leanbox}[1]{enhanced,breakable,colback=leanback,colframe=leanrule,
  boxrule=0.4pt,arc=0.6mm,left=2.4mm,right=2.4mm,top=1.8mm,bottom=1.8mm,
  before skip=0.8em,after skip=1em,fonttitle=\small\bfseries,coltitle=leantitle,
  colbacktitle=leanback,titlerule=0.3pt,title={#1}}
\newlength{\leancharwidth}
\newenvironment{leancode}{\par\smallskip\begingroup\ttfamily\footnotesize\raggedright
  \settowidth{\leancharwidth}{M}\setlength{\parskip}{0pt}\setlength{\parindent}{0pt}
  \frenchspacing}{\par\endgroup\smallskip}
\newcommand{\leanline}[2]{\hangindent=\dimexpr#1\leancharwidth+4\leancharwidth\relax
  \hangafter=1\noindent\hspace*{#1\leancharwidth}#2\par}
\newcommand{\leansource}[1]{\par{\footnotesize\color{leanmuted}#1}\par}
\newcommand{\leanlabel}[1]{\par\smallskip\noindent{\small\bfseries #1}\par}
\usepackage[margin=25mm]{geometry}
\usepackage[unicode,hidelinks]{hyperref}
\hypersetup{pdfauthor={LangevinLean Project}}
\setlength{\emergencystretch}{3em}
\setlength{\parskip}{0.4em}
\setlength{\parindent}{0pt}
\allowdisplaybreaks[2]
\raggedbottom
\title{\texorpdfstring{Kinetic Langevin Monte Carlo}{Kinetic Langevin Monte Carlo}}
\author{LangevinLean Project}
\date{}
\begin{document}\maketitle
Draft mathematical exposition; these notes may contain mistakes.
The Brownian kinetic mixing theorem, first-order KLMC sampler theorem, Hessian-corrected KLMC2 theorem, Gibbs stationarity, exact Gaussian and Cholesky constructions, and explicit accuracy corollaries are checked. The source audit records the supplied manuscript corrections and the continuous-\(L^2\) Brownian model.
\section{\texorpdfstring{Theorem}{Theorem}}


Let \(E\) be a complete real inner product space with its Borel
structure and a countable topological base. Suppose \(f:E\to\mathbb R\) is
differentiable and \(m\)-strongly convex with \(M\)-Lipschitz gradient,
\(m\le M\). Fix \(\gamma\ge0\), \(a<\gamma/2\), and \(0\le t\le T\).
Let \(\mu,\nu\) be probability laws of the initial position with
\(W_2(\mu,\nu)<\infty\). Let \(Q\) be any probability law of the initial
velocity and continuous velocity forcing on \([0,T]\), independent of the
initial position in each process.




The selected flow solves the integral equation, with normalized
forcing \(b_t-b_0\),


\[Y_t-Y_0=-\int_0^t(\gamma V_s+\nabla f(X_s),-V_s)ds+(b_t-b_0,0).\]


Denote its position law by \(P_t^Q\mu\). Then


\[\boxed{W_2(P_t^Q\mu,P_t^Q\nu)\le A_a e^{k_at}W_2(\mu,\nu),}\]
\[A_a=\frac{\sqrt{2((\gamma-a)^2+a^2)}}{\gamma-2a},\]
\[k_a=\frac{\max\{a^2-m,M-(\gamma-a)^2\}}{\gamma-2a}.\]


This is the explicit prefactor and exponent of the manuscript's equation
(exp2), proved for the independent-product laws and the infimum over
initial couplings. No continuous Hessian assumption is needed. Choosing
\(a=0\) and \(\gamma^2\ge m+M\) gives
\(\sqrt2e^{-mt/\gamma}W_2(\mu,\nu)\).




The proof first cancels the common forcing in two integral solutions.
The earlier strong-smooth gradient interpolation theorem (see \cite[Section~3.10.2]{LangevinLean2026ULASC}) controls the transformed energy
\(\|V-V'+(\gamma-a)(X-X')\|^2+\|-(V-V')-a(X-X')\|^2\).
Gronwall gives the path estimate. Each initial coupling is then extended by
an independent copy of \(Q\); both output marginals and its cost are proved
before applying the earlier all-couplings transport infimum theorem (see \cite[Section~3.18.4]{LangevinLean2026ULASC}). The specified Brownian source and
its Gibbs-stationary solution are also constructed in the supporting modules;
the separate KLMC theorem below combines them with the numerical analysis.



The source law may include any common square-integrable velocity and admissible continuous forcing, provided it is independent of each initial position. The noise distribution does not enter the difference estimate because the same realization drives both trajectories. The result therefore separates deterministic contraction from the later specialization to Brownian dynamics.

\begin{leanbox}{Lean: \texttt{kineticPositionLaw\_\allowbreak{}wasserstein2Real\_\allowbreak{}general}}
\leansource{\texttt{kineticPositionLaw\_\allowbreak{}wasserstein2Real\_\allowbreak{}general} \quad KineticLangevin/\allowbreak{}FlowLaw.lean, lines 270–285}
\begin{leancode}
\leanline{0}{theorem\ kineticPositionLaw\_wasserstein2Real\_general\ (f\ :\ E\ →\ ℝ)\ (m\ a\ :\ ℝ)\ (M\ γ\ :\ ℝ≥0)}
\leanline{4}{(hdiff\ :\ Differentiable\ ℝ\ f)\ (hsc\ :\ StrongConvexOn\ univ\ m\ f)}
\leanline{4}{(hLip\ :\ LipschitzWith\ M\ (gradient\ f))\ (hmM\ :\ m\ ≤\ (M\ :\ ℝ))}
\leanline{4}{(hd\ :\ 0\ <\ (γ\ :\ ℝ)-2*a)}
\leanline{4}{(T\ :\ ℝ)\ (hT\ :\ 0\ ≤\ T)\ (t\ :\ ℝ)\ (ht\ :\ t\ ∈\ Icc\ 0\ T)}
\leanline{4}{(μ\ ν\ :\ Measure\ E)\ [IsProbabilityMeasure\ μ]\ [IsProbabilityMeasure\ ν]}
\leanline{4}{(Q\ :\ Measure\ (E\ ×\ VelocityForcing\ (E\ :=\ E)\ T))\ [IsProbabilityMeasure\ Q]}
\leanline{4}{(hfinite\ :\ ULA.wasserstein2\ μ\ ν\ ≠\ ⊤)\ :}
\leanline{4}{ULA.wasserstein2Real\ (kineticPositionLaw\ (gradient\ f)\ M\ γ\ hLip\ T\ hT\ t\ μ\ Q)}
\leanline{6}{(kineticPositionLaw\ (gradient\ f)\ M\ γ\ hLip\ T\ hT\ t\ ν\ Q)\ ≤}
\leanline{6}{(Real.sqrt\ (2\ *\ (((γ\ :\ ℝ)-a)\textasciicircum{}2+a\textasciicircum{}2))\ /\ ((γ\ :\ ℝ)-2*a)\ *}
\leanline{8}{Real.exp\ ((max\ (a\textasciicircum{}2-m)\ ((M\ :\ ℝ)-((γ\ :\ ℝ)-a)\textasciicircum{}2)\ /\ ((γ\ :\ ℝ)-2*a))*t))\ *}
\leanline{8}{ULA.wasserstein2Real\ μ\ ν}
\end{leancode}
\end{leanbox}

\section{\texorpdfstring{The kinetic equation and its Gibbs phase measure}{The kinetic equation and its Gibbs phase measure}}


Let the target potential be \(f\): the target density is proportional to \(e^{-f}\). The kinetic process has a velocity variable and a position variable satisfying


\[dV_t=-(\gamma V_t+\nabla f(L_t))dt+\sqrt{2\gamma}\,dW_t,
\qquad dL_t=V_tdt.\]


The linear friction \(\gamma V_t\) removes kinetic energy. The noise amplitude gives a standard Gaussian equilibrium velocity.
The candidate equilibrium has density


\[\Pi(dx,dv)=\pi(dx)\,\varphi(dv)
\propto e^{-f(x)-\|v\|^2/2}dx\,dv.\]


Position and velocity are independent under this law. Integrating out velocity
therefore leaves precisely the desired target. Thus the position marginal of this phase measure is the target sampling distribution.




Strong convexity supplies a quadratic lower bound on \(f\), so the partition
function is finite and the required position moments exist (see \cite[Section~3.26.10]{LangevinLean2026ULASC}, \cite[Section~3.26.15]{LangevinLean2026ULASC}). A proof of
stationarity still needs a solution process and a valid invariance
argument; checking that a formal generator integrates to zero is one component
of that argument.



The kinetic extension is useful because position is the time integral of velocity even though velocity contains Brownian forcing. This extra time integration reduces the local positional discretization defect. At equilibrium, the Gaussian velocity radius is exactly \(\sqrt p\) in second mean, which determines the dimension factor in the first-order error theorem.

\begin{leanbox}{Lean: \texttt{phaseFlow\_\allowbreak{}gibbs\_\allowbreak{}stationary}, \texttt{positionLaw\_\allowbreak{}gibbs}}
\leansource{\texttt{phaseFlow\_\allowbreak{}gibbs\_\allowbreak{}stationary} \quad KineticLangevin/\allowbreak{}ContinuousGibbs.lean, lines 89–118}
\begin{leancode}
\leanline{0}{theorem\ phaseFlow\_gibbs\_stationary\ (f\ :\ H\ ι\ →\ ℝ)\ (m\ :\ ℝ)\ (hm\ :\ 0\ <\ m)\ (M\ γ\ :\ ℝ≥0)}
\leanline{4}{(hf\ :\ Differentiable\ ℝ\ f)\ (hsc\ :\ StrongConvexOn\ univ\ m\ f)\ (hLip\ :\ LipschitzWith\ M\ (gradient\ f))}
\leanline{4}{(T\ :\ ℝ)\ (hT\ :\ 0\ ≤\ T)\ :}
\leanline{4}{letI\ :=\ ULA.isProbabilityMeasure\_gibbsMeasure\ f\ m\ hm\ hf\ hsc}
\leanline{4}{HasLaw\ (lpPhaseCoe\ ((ULA.gibbsMeasure\ f).prod\ (probability\ ι\ (stdGaussian\ (H\ ι))))}
\leanline{6}{(ProductContinuous.flow\ (probability\ ι\ (stdGaussian\ (H\ ι)))\ (ULA.gibbsMeasure\ f)}
\leanline{8}{(ULA.memLp\_id\_gibbsMeasure\ f\ m\ hm\ hf\ hsc)\ (velocity\ ι\ (stdGaussian\ (H\ ι))\ IsGaussian.memLp\_two\_id)}
\leanline{8}{(input\ ι\ (stdGaussian\ (H\ ι))\ γ)\ (continuous\_input\ ι\ (stdGaussian\ (H\ ι))\ γ)}
\leanline{8}{(gradient\ f)\ M\ γ\ hLip\ T\ hT\ T))}
\leanline{6}{(kineticGibbsTarget\ f)\ ((ULA.gibbsMeasure\ f).prod\ (probability\ ι\ (stdGaussian\ (H\ ι))))}
\end{leancode}
\leansource{\texttt{positionLaw\_\allowbreak{}gibbs} \quad KineticLangevin/\allowbreak{}ContinuousGibbs.lean, lines 121–130}
\begin{leancode}
\leanline{0}{theorem\ positionLaw\_gibbs\ (f\ :\ H\ ι\ →\ ℝ)\ (m\ :\ ℝ)\ (hm\ :\ 0\ <\ m)\ (M\ γ\ :\ ℝ≥0)}
\leanline{4}{(hf\ :\ Differentiable\ ℝ\ f)\ (hsc\ :\ StrongConvexOn\ univ\ m\ f)\ (hLip\ :\ LipschitzWith\ M\ (gradient\ f))}
\leanline{4}{(T\ :\ ℝ)\ (hT\ :\ 0\ ≤\ T)\ :}
\leanline{4}{letI\ :=\ ULA.isProbabilityMeasure\_gibbsMeasure\ f\ m\ hm\ hf\ hsc}
\leanline{4}{positionLaw\ ι\ (stdGaussian\ (H\ ι))\ IsGaussian.memLp\_two\_id\ (ULA.gibbsMeasure\ f)}
\leanline{6}{(ULA.memLp\_id\_gibbsMeasure\ f\ m\ hm\ hf\ hsc)\ f\ M\ γ\ hLip\ T\ hT\ =\ ULA.gibbsMeasure\ f}
\end{leancode}
\leanlabel{How the Lean proof follows the argument}
\texttt{phase\allowbreak{}Flow\_\allowbreak{}gibbs\_\allowbreak{}stationary} first installs the normalized Gibbs probability instance and its second-moment bound. The local objects \texttt{Q}, \texttt{V}, \texttt{B} and \texttt{Y} specify the independent Gaussian/\allowbreak{}Brownian source, initial velocity, forcing and integral solution. \texttt{hs} supplies the integral equation and continuity; \texttt{hjoint} supplies the complete initial-phase/\allowbreak{}OU-increment law. These are the inputs of \texttt{kinetic\_\allowbreak{}integral\_\allowbreak{}stationary\_\allowbreak{}of\_\allowbreak{}joint\allowbreak{}Noise}. Thus invariance is obtained for the constructed solution. \texttt{position\allowbreak{}Law\_\allowbreak{}gibbs} composes the resulting \texttt{Has\allowbreak{}Law} with the second-coordinate projection to recover the Gibbs position marginal.
\end{leanbox}

\section{\texorpdfstring{Normalization of inverse mass}{Normalization of inverse mass}}


The general paper equation includes inverse mass \(u>0\):
\(dV=-(\gamma V+u\nabla f(L))dt+\sqrt{2\gamma u}\,dW\).
Its equilibrium velocity has covariance \(uI\). Normalize both the velocity
and the time variable:


\[\bar V_t=u^{-1/2}V_{t/\sqrt u},\qquad
\bar L_t=L_{t/\sqrt u},\qquad \bar\gamma=\gamma/\sqrt u.\]


The position equation becomes \(d\bar L_t=\bar V_tdt\). Brownian scaling
turns \(u^{1/4}W_{t/\sqrt u}\) into a standard Brownian motion; substitution
in the velocity equation gives the same kinetic system with \(u=1\) and
friction \(\bar\gamma\). This is Lemma 1 in the supplied manuscript.




The scaling explains why the rest of the paper fixes \(u=1\): it removes
a redundant clock parameter. When transferring a theorem back to another mass,
the step size, friction, Gaussian velocity law and physical time must all be
transformed together.



To check the transformation, write the position increment as an integral before changing variables: \(L_{t/\sqrt u}-L_0=\int_0^{t/\sqrt u}V_sds=\int_0^tu^{-1/2}V_{r/\sqrt u}dr\). The velocity equation transforms similarly. This verifies the normalized equation for integral solutions, including its initial data, rather than only manipulating differential notation.

\begin{leanbox}{Lean: \texttt{kinetic\_\allowbreak{}mass\_\allowbreak{}integral}, \texttt{kinetic\_\allowbreak{}mass\_\allowbreak{}noise\_\allowbreak{}amplitude}, \texttt{clockBrownian\_\allowbreak{}isPreBrownian}}
\leansource{\texttt{kinetic\_\allowbreak{}mass\_\allowbreak{}integral} \quad KineticLangevin/\allowbreak{}MassScaling.lean, lines 28–54}
\begin{leancode}
\leanline{0}{theorem\ kinetic\_mass\_integral\ (g\ :\ E→E)\ (hg\ :\ Continuous\ g)\ (c\ γ\ T\ :\ ℝ)\ (hc\ :\ 0<c)}
\leanline{4}{(initial\ :\ E×E)\ (B\ :\ ℝ→E)\ (Y\ :\ ℝ→E×E)\ (hY\ :\ Continuous\ Y)}
\leanline{4}{(heY\ :\ ∀\ s\ ∈\ Icc\ 0\ (T/c),Y\ s-initial\ =}
\leanline{6}{-(∫\ r\ in\ 0..s,kineticNegativeDrift\ g\ γ\ (c\textasciicircum{}2)\ (Y\ r))+(B\ s,0))\ :}
\leanline{4}{∀\ t\ ∈\ Icc\ 0\ T,}
\leanline{6}{massPhaseMap\ c⁻¹\ (Y\ (t/c))-massPhaseMap\ c⁻¹\ initial\ =}
\leanline{8}{-(∫\ s\ in\ 0..t,kineticNegativeDrift\ g\ (γ/c)\ 1\ (massPhaseMap\ c⁻¹\ (Y\ (s/c))))+}
\leanline{10}{(c⁻¹\ •\ B\ (t/c),0)}
\end{leancode}
\leansource{\texttt{kinetic\_\allowbreak{}mass\_\allowbreak{}noise\_\allowbreak{}amplitude} \quad KineticLangevin/\allowbreak{}MassScaling.lean, lines 57–62}
\begin{leancode}
\leanline{0}{theorem\ kinetic\_mass\_noise\_amplitude\ (c\ γ\ :\ ℝ)\ (hc\ :\ 0<c)\ (hγ\ :\ 0≤γ)\ :}
\leanline{4}{c⁻¹*Real.sqrt\ (2*γ*c\textasciicircum{}2)=Real.sqrt\ (2*(γ/c))*Real.sqrt\ c}
\end{leancode}
\leansource{\texttt{clockBrownian\_\allowbreak{}isPreBrownian} \quad KineticLangevin/\allowbreak{}BrownianScaling.lean, lines 15–19}
\begin{leancode}
\leanline{0}{theorem\ clockBrownian\_isPreBrownian\ \{ι\ :\ Type*\}\ (c\ :\ ℝ≥0)\ (hc\ :\ c≠0)\ (i\ :\ ι)\ :}
\leanline{4}{IsPreBrownianReal\ (clockBrownian\ c\ i)\ (ULA.CanonicalBrownian.vectorProbability\ ι)}
\end{leancode}
\leanlabel{How the Lean proof follows the argument}
\texttt{kinetic\_\allowbreak{}mass\_\allowbreak{}integral} checks the velocity and position integral equations after the deterministic clock and velocity transformations. \texttt{kinetic\_\allowbreak{}mass\_\allowbreak{}noise\_\allowbreak{}amplitude} checks the resulting square-root coefficient. \texttt{clock\allowbreak{}Brownian\_\allowbreak{}is\allowbreak{}Pre\allowbreak{}Brownian} separately verifies the Brownian law at the changed clock; it supplies the probabilistic justification that cannot be obtained by algebraically substituting time in a stochastic differential notation.
\end{leanbox}

\section{\texorpdfstring{The difference equation under common Brownian forcing}{The difference equation under common Brownian forcing}}


Consider two kinetic solutions with different initial positions and the same Brownian input. The random forcing cancels in their difference. Put
\(q_t=L_t-L'_t\), \(r_t=V_t-V'_t\). Then


\[q'_t=r_t,\qquad r'_t=-\gamma r_t-H_tq_t,\qquad
H_t=\int_0^1\nabla^2f(L'_t+s(L_t-L'_t))\,ds.\]


The averaged Hessian is symmetric and obeys \(mI\preceq H_t\preceq MI\).
The problem has become a deterministic differential inequality along each pair
of paths. Existence and uniqueness of those common-input paths follow by applying the earlier additive-input solution theorems to the globally Lipschitz kinetic phase drift (see \cite[Section~3.14.3]{LangevinLean2026ULASC}, \cite[Section~3.43.2]{LangevinLean2026ULASC}). No independent-noise variance is added to this comparison.




The usual Euclidean norm of \((r,q)\) need not decrease monotonically:
energy can move between position and velocity. The paper instead chooses a
linear transformation whose squared norm satisfies a differential inequality.



The cancellation takes place before expectations are taken. Consequently the same deterministic estimate holds on every sample path where the integral equations hold, and also in the Hilbert space of square-integrable random vectors. Independent forcing would instead introduce an extra random difference and would not produce this contraction comparison.

\begin{leanbox}{Lean: \texttt{kinetic\_\allowbreak{}integral\_\allowbreak{}position\_\allowbreak{}decay}, \texttt{wasserstein2\_\allowbreak{}gibbs\_\allowbreak{}general}}
\leansource{\texttt{kinetic\_\allowbreak{}integral\_\allowbreak{}position\_\allowbreak{}decay} \quad KineticLangevin/\allowbreak{}SynchronousFlow.lean, lines 64–94}
\begin{leancode}
\leanline{0}{theorem\ kinetic\_integral\_position\_decay}
\leanline{4}{(f\ :\ E\ →\ ℝ)\ (m\ :\ ℝ)\ (M\ γ\ :\ ℝ≥0)}
\leanline{4}{(hdiff\ :\ Differentiable\ ℝ\ f)\ (hsc\ :\ StrongConvexOn\ univ\ m\ f)}
\leanline{4}{(hLip\ :\ LipschitzWith\ M\ (gradient\ f))\ (hmM\ :\ m\ ≤\ (M\ :\ ℝ))}
\leanline{4}{(hγ\ :\ 0\ <\ (γ\ :\ ℝ))\ (hfriction\ :\ m\ +\ (M\ :\ ℝ)\ ≤\ (γ\ :\ ℝ)\ \textasciicircum{}\ 2)}
\leanline{4}{(Y\ Z\ B\ :\ ℝ\ →\ E\ ×\ E)\ (hY\ :\ Continuous\ Y)\ (hZ\ :\ Continuous\ Z)\ (T\ :\ ℝ)}
\leanline{4}{(heY\ :\ ∀\ t\ ∈\ Icc\ 0\ T,\ Y\ t\ -\ Y\ 0\ =}
\leanline{6}{-(∫\ s\ in\ 0..t,\ kineticNegativeDrift\ (gradient\ f)\ γ\ 1\ (Y\ s))\ +\ B\ t)}
\leanline{4}{(heZ\ :\ ∀\ t\ ∈\ Icc\ 0\ T,\ Z\ t\ -\ Z\ 0\ =}
\leanline{6}{-(∫\ s\ in\ 0..t,\ kineticNegativeDrift\ (gradient\ f)\ γ\ 1\ (Z\ s))\ +\ B\ t)}
\leanline{4}{(hvelocity\ :\ (Y\ 0).1\ =\ (Z\ 0).1)\ (t\ :\ ℝ)\ (ht\ :\ t\ ∈\ Icc\ 0\ T)\ :}
\leanline{4}{‖(Y\ t).2\ -\ (Z\ t).2‖\ ≤\ Real.sqrt\ 2\ *\ Real.exp\ ((-m\ /\ γ)\ *\ t)\ *}
\leanline{6}{‖(Y\ 0).2\ -\ (Z\ 0).2‖}
\end{leancode}
\leansource{\texttt{wasserstein2\_\allowbreak{}gibbs\_\allowbreak{}general} \quad KineticLangevin/\allowbreak{}ContinuousGibbs.lean, lines 133–147}
\begin{leancode}
\leanline{0}{theorem\ wasserstein2\_gibbs\_general\ (μ\ :\ Measure\ (H\ ι))\ [IsProbabilityMeasure\ μ]}
\leanline{4}{(hμ\ :\ MemLp\ (id\ :\ H\ ι\ →\ H\ ι)\ 2\ μ)\ (f\ :\ H\ ι\ →\ ℝ)\ (m\ a\ :\ ℝ)\ (M\ γ\ :\ ℝ≥0)}
\leanline{4}{(hm\ :\ 0\ <\ m)\ (hf\ :\ Differentiable\ ℝ\ f)\ (hsc\ :\ StrongConvexOn\ univ\ m\ f)}
\leanline{4}{(hLip\ :\ LipschitzWith\ M\ (gradient\ f))\ (hmM\ :\ m\ ≤\ (M\ :\ ℝ))\ (hd\ :\ 0\ <\ (γ\ :\ ℝ)-2*a)}
\leanline{4}{(T\ :\ ℝ)\ (hT\ :\ 0\ ≤\ T)\ :}
\leanline{4}{ULA.wasserstein2Real\ (positionLaw\ ι\ (stdGaussian\ (H\ ι))\ IsGaussian.memLp\_two\_id\ μ\ hμ\ f\ M\ γ\ hLip\ T\ hT)}
\leanline{6}{(ULA.gibbsMeasure\ f)\ ≤}
\leanline{6}{(Real.sqrt\ (2*(((γ\ :\ ℝ)-a)\textasciicircum{}2+a\textasciicircum{}2))/((γ\ :\ ℝ)-2*a)\ *}
\leanline{8}{Real.exp\ (max\ (a\textasciicircum{}2-m)\ ((M\ :\ ℝ)-((γ\ :\ ℝ)-a)\textasciicircum{}2)/((γ\ :\ ℝ)-2*a)*T))\ *}
\leanline{10}{ULA.wasserstein2Real\ μ\ (ULA.gibbsMeasure\ f)}
\end{leancode}
\leanlabel{How the Lean proof follows the argument}
\texttt{kinetic\_\allowbreak{}integral\_\allowbreak{}position\_\allowbreak{}decay} subtracts the two common-input integral equations, so the additive forcing disappears before differentiation and the transformed energy inequality is used. \texttt{wasserstein2\_\allowbreak{}gibbs\_\allowbreak{}general} applies the corresponding two-initial-law coupling theorem and rewrites the stationary comparator using \texttt{position\allowbreak{}Law\_\allowbreak{}gibbs}. The initial coupling is arbitrary; the transport theorem passes its cost to the infimum defining Wasserstein distance.
\end{leanbox}

\section{\texorpdfstring{The transformed energy and its derivative}{The transformed energy and its derivative}}


Choose \(0\le v<\gamma/2\), set \(\lambda_-=v\),
\(\lambda_+=\gamma-v\), and define


\[a_t=r_t+\lambda_+q_t,\qquad b_t=-r_t-\lambda_-q_t,
\qquad \mathcal E_t=\|a_t\|^2+\|b_t\|^2.\]


The change of variables is invertible because \(\lambda_+-\lambda_->0\).
Differentiate \(\mathcal E\), use the equations for \(q,r\), and exploit
symmetry of \(H_t\). The two off-diagonal inner products cancel exactly:


\[\mathcal E'_t=
\frac{2}{\lambda_+-\lambda_-}
\bigl(\langle a_t,(\lambda_-^2I-H_t)a_t\rangle
+\langle b_t,(H_t-\lambda_+^2I)b_t\rangle\bigr).\]


Consequently \(\mathcal E'_t\le2c(v)\mathcal E_t\), where


\[c(v)=\frac{\max\{v^2-m,\ M-(\gamma-v)^2\}}{\gamma-2v}.\]


Gronwall gives \(\sqrt{\mathcal E_t}\le e^{c(v)t}\sqrt{\mathcal E_0}\).
If the initial velocities agree, then
\(\mathcal E_0=((\gamma-v)^2+v^2)\|q_0\|^2\). Since
\(q_t=(a_t+b_t)/(\gamma-2v)\), the position estimate becomes


\[\|q_t\|\le
\frac{\sqrt{2((\gamma-v)^2+v^2)}}{\gamma-2v}
e^{c(v)t}\|q_0\|.\]


This is the explicit form of Proposition 1 and Theorem 1. Squaring and integrating first bounds the output cost for each initial coupling. The earlier coupling-cost bound and all-couplings infimum theorem then turn these bounds into a Wasserstein estimate (see \cite[Section~3.17.13]{LangevinLean2026ULASC}, \cite[Section~3.18.4]{LangevinLean2026ULASC}).



The inverse position identity follows by adding the two transformed coordinates. Its norm bound is \(\|a_t+b_t\|^2\le2(\|a_t\|^2+\|b_t\|^2)\), obtained from the triangle inequality and \(2rs\le r^2+s^2\). This is the sole projection loss. Propagating the transformed norm through the entire evolution and projecting only once keeps that projection coefficient outside the iteration.

\begin{leanbox}{Lean: \texttt{transform}, \texttt{spectral\_\allowbreak{}identity}, \texttt{potential\_\allowbreak{}general\_\allowbreak{}force\_\allowbreak{}dissipativity}}
\leansource{\texttt{transform} \quad KineticLangevin/\allowbreak{}Geometry.lean, lines 20–21}
\begin{leancode}
\leanline{0}{def\ transform\ (γ\ a\ :\ ℝ)\ (s\ :\ E\ ×\ E)\ :\ E\ ×\ E\ :=}
\leanline{2}{(s.1\ +\ (γ\ -\ a)\ •\ s.2,\ -s.1\ -\ a\ •\ s.2)}
\end{leancode}
\leansource{\texttt{spectral\_\allowbreak{}identity} \quad KineticLangevin/\allowbreak{}Contraction.lean, lines 28–40}
\begin{leancode}
\leanline{0}{theorem\ spectral\_identity\ (γ\ a\ :\ ℝ)\ (H\ :\ E\ →L[ℝ]\ E)}
\leanline{4}{(hH\ :\ ∀\ x\ y,\ {\leanfallback ⟪}x,\ H\ y{\leanfallback ⟫}\ =\ {\leanfallback ⟪}y,\ H\ x{\leanfallback ⟫})\ (s\ :\ E\ ×\ E)\ :}
\leanline{4}{(γ\ -\ 2\ *\ a)\ *\ energyDerivative\ γ\ a\ H\ s\ =}
\leanline{6}{2\ *\ (a\ \textasciicircum{}\ 2\ *\ ‖(transform\ γ\ a\ s).1‖\ \textasciicircum{}\ 2\ -}
\leanline{8}{{\leanfallback ⟪}(transform\ γ\ a\ s).1,\ H\ (transform\ γ\ a\ s).1{\leanfallback ⟫}\ +}
\leanline{8}{{\leanfallback ⟪}(transform\ γ\ a\ s).2,\ H\ (transform\ γ\ a\ s).2{\leanfallback ⟫}\ -}
\leanline{8}{(γ\ -\ a)\ \textasciicircum{}\ 2\ *\ ‖(transform\ γ\ a\ s).2‖\ \textasciicircum{}\ 2)}
\end{leancode}
\leansource{\texttt{potential\_\allowbreak{}general\_\allowbreak{}force\_\allowbreak{}dissipativity} \quad KineticLangevin/\allowbreak{}GeneralContraction.lean, lines 116–126}
\begin{leancode}
\leanline{0}{theorem\ potential\_general\_force\_dissipativity\ [CompleteSpace\ E]}
\leanline{4}{(f\ :\ E\ →\ ℝ)\ (m\ γ\ a\ :\ ℝ)\ (M\ :\ ℝ≥0)}
\leanline{4}{(hdiff\ :\ Differentiable\ ℝ\ f)\ (hsc\ :\ StrongConvexOn\ univ\ m\ f)}
\leanline{4}{(hLip\ :\ LipschitzWith\ M\ (gradient\ f))\ (hmM\ :\ m\ ≤\ (M\ :\ ℝ))}
\leanline{4}{(hd\ :\ 0\ <\ γ\ -\ 2\ *\ a)\ (v\ x\ y\ :\ E)\ :}
\leanline{4}{generalForceEnergyDerivative\ γ\ a\ v\ (x\ -\ y)\ (gradient\ f\ x\ -\ gradient\ f\ y)\ ≤}
\leanline{6}{(2\ *\ max\ (a\ \textasciicircum{}\ 2\ -\ m)\ ((M\ :\ ℝ)\ -\ (γ\ -\ a)\ \textasciicircum{}\ 2)\ /\ (γ\ -\ 2\ *\ a))\ *}
\leanline{8}{transformEnergy\ γ\ a\ (v,\ x\ -\ y)}
\end{leancode}
\leanlabel{How the Lean proof follows the argument}
\texttt{transform} is exactly the two-coordinate linear change of variables in the argument; its energy is a sum of squared norms, rather than the default maximum norm of a product type. \texttt{spectral\_\allowbreak{}identity} expands the inner products in the quadratic dissipation. \texttt{potential\_\allowbreak{}general\_\allowbreak{}force\_\allowbreak{}dissipativity} uses strong convexity and gradient Lipschitz continuity to bound that quadratic form for the actual force difference. The two scalar spectral bounds correspond to the two entries of the maximum in the displayed exponent.
\end{leanbox}

\section{\texorpdfstring{The friction regimes and their contraction coefficients}{The friction regimes and their contraction coefficients}}


At \(v=0\) the transformed energy is particularly simple:
\(\|r+\gamma q\|^2+\|r\|^2\). If \(\gamma^2\ge m+M\), then
\(c(0)=-m/\gamma\), so


\[W_2(\mu P_t^L,\mu'P_t^L)
\le\sqrt2e^{-mt/\gamma}W_2(\mu,\mu').\]


Here the two initial phase-space laws share the same independent velocity
marginal. The prefactor is \(\sqrt2\); this follows directly by substituting
\(v=0\) in the explicit formula. The supplied source's earlier
\(\sqrt2/\gamma\) prefactor in its unshifted display does not match
that substitution.




Friction that is very large makes the rate \(m/\gamma\) small: velocities
are damped before they transport position very far. The smallest friction
allowed by this simple estimate is \(\sqrt{m+M}\). The full table in the
paper improves the rate in other regimes by optimizing \(v\). For sufficiently
large friction it recovers


\[\beta=\frac{\gamma-\sqrt{\gamma^2-4m}}2,\]


This is the slow characteristic rate of a quadratic potential with curvature \(m\).
This explains the sharpness example. Singular prefactors at a repeated
characteristic root require separate limiting estimates.



The optimized exponent and the prefactor describe different effects. A more negative exponent improves long-time decay, while a large prefactor can make a finite-time bound less useful. Near the repeated-root boundary the coordinate transformation becomes singular, so rate optimization must retain its domain and cannot simply substitute the boundary into the prefactor.

\begin{leanbox}{Lean: \texttt{contractionExponent\_\allowbreak{}low\_\allowbreak{}optimal}, \texttt{contractionExponent\_\allowbreak{}middle\_\allowbreak{}optimal}, \texttt{contractionExponent\_\allowbreak{}high\_\allowbreak{}optimal}}
\leansource{\texttt{contractionExponent\_\allowbreak{}low\_\allowbreak{}optimal} \quad KineticLangevin/\allowbreak{}Friction.lean, lines 37–44}
\begin{leancode}
\leanline{0}{theorem\ contractionExponent\_low\_optimal\ (γ\ m\ M\ a\ :\ ℝ)\ (hγ\ :\ 0\ <\ γ)}
\leanline{4}{(hmM\ :\ m\ ≤\ M)\ (hregime\ :\ γ\textasciicircum{}2\ ≤\ m+M)\ (ha\ :\ 0\ ≤\ a)\ (hd\ :\ 0\ <\ γ-2*a)\ :}
\leanline{4}{(M-γ\textasciicircum{}2)/γ\ ≤\ contractionExponent\ γ\ a\ m\ M}
\end{leancode}
\leansource{\texttt{contractionExponent\_\allowbreak{}middle\_\allowbreak{}optimal} \quad KineticLangevin/\allowbreak{}Friction.lean, lines 77–107}
\begin{leancode}
\leanline{0}{theorem\ contractionExponent\_middle\_optimal\ (γ\ m\ M\ a\ :\ ℝ)\ (hγ\ :\ 0\ <\ γ)}
\leanline{4}{(hmM\ :\ m\ ≤\ M)\ (hleft\ :\ m+M\ ≤\ γ\textasciicircum{}2)\ (hright\ :\ γ\textasciicircum{}2\ ≤\ 3*m+M)}
\leanline{4}{(hradicand\ :\ 0\ <\ 2*(m+M)-γ\textasciicircum{}2)\ (ha\ :\ 0\ ≤\ a)\ (hd\ :\ 0\ <\ γ-2*a)\ :}
\leanline{4}{-γ/2\ +\ (M-m)/(2*Real.sqrt\ (2*(m+M)-γ\textasciicircum{}2))\ ≤\ contractionExponent\ γ\ a\ m\ M}
\end{leancode}
\leansource{\texttt{contractionExponent\_\allowbreak{}high\_\allowbreak{}optimal} \quad KineticLangevin/\allowbreak{}Friction.lean, lines 131–139}
\begin{leancode}
\leanline{0}{theorem\ contractionExponent\_high\_optimal\ (γ\ m\ M\ a\ :\ ℝ)\ (hdisc\ :\ 0\ <\ γ\textasciicircum{}2-4*m)}
\leanline{4}{(hd\ :\ 0\ <\ γ-2*a)\ :}
\leanline{4}{-(γ-Real.sqrt\ (γ\textasciicircum{}2-4*m))/2\ ≤\ contractionExponent\ γ\ a\ m\ M}
\end{leancode}
\leanlabel{How the Lean proof follows the argument}
\texttt{contraction\allowbreak{}Exponent\_\allowbreak{}low\_\allowbreak{}optimal}, \texttt{contraction\allowbreak{}Exponent\_\allowbreak{}middle\_\allowbreak{}optimal} and \texttt{contraction\allowbreak{}Exponent\_\allowbreak{}high\_\allowbreak{}optimal} are the three scalar comparisons for the admissible shift. Their case hypotheses select the active branch of the maximum. Each conclusion compares the candidate exponent with every admissible shift; the nonsingularity condition keeps the transformation denominator positive. The full shifted Brownian theorem is retained when a simple-friction specialization is unavailable.
\end{leanbox}

\section{\texorpdfstring{The exponential update with frozen force}{The exponential update with frozen force}}


On an interval of length \(h\), replace the changing force
\(\nabla f(L_t)\) by its value at the current position \(x\). The remaining
velocity equation is linear and can be solved exactly. Define


\[\psi_0(h)=e^{-\gamma h},\qquad
\psi_1(h)=\frac{1-e^{-\gamma h}}\gamma,\qquad
\psi_2(h)=\frac h\gamma-\frac{1-e^{-\gamma h}}{\gamma^2}.\]


With velocity written first, the update is


\[v^+=\psi_0(h)v-\psi_1(h)\nabla f(x)+\sqrt{2\gamma}\,\xi,
\quad x^+=x+\psi_1(h)v-\psi_2(h)\nabla f(x)+\sqrt{2\gamma}\,\xi'.\]


The exponential terms handle friction at its exact continuous-time rate.
The coefficients of the force are its first and second time integrals. For
small \(h\), \(\psi_1(h)\approx h\) and \(\psi_2(h)\approx h^2/2\),
which are the first-order velocity and second-order position changes
under a constant force.




The only approximation here is freezing the nonlinear force. The linear
friction and its associated Gaussian forcing have already been integrated.
This is the KLMC exponential integrator reused by KIPLMC1.



An integrating factor gives \(v_t=e^{-\gamma t}v-\int_0^te^{-\gamma(t-s)}\nabla f(x)ds\) before noise is added. Since the force is constant inside this auxiliary problem, the integral equals \(\psi_1(t)\nabla f(x)\). Integrating this velocity once gives the \(\psi_2\) position coefficient. The two coefficient formulas therefore arise from the same exact frozen system.

\begin{leanbox}{Lean: \texttt{exponentialStep}, \texttt{klmcStep}}
\leansource{\texttt{exponentialStep} \quad KineticLangevin/\allowbreak{}Integrator.lean, lines 57–60}
\begin{leancode}
\leanline{0}{def\ exponentialStep\ (force\ :\ E\ →\ E)\ (γ\ h\ :\ ℝ)}
\leanline{4}{(state\ noise\ :\ PhaseSpace\ E)\ :\ PhaseSpace\ E\ :=}
\leanline{2}{(psi\ γ\ 0\ h\ •\ state.1\ -\ psi\ γ\ 1\ h\ •\ force\ state.2\ +\ noise.1,}
\leanline{3}{state.2\ +\ psi\ γ\ 1\ h\ •\ state.1\ -\ psi\ γ\ 2\ h\ •\ force\ state.2\ +\ noise.2)}
\end{leancode}
\leansource{\texttt{klmcStep} \quad KineticLangevin/\allowbreak{}Integrator.lean, lines 70–72}
\begin{leancode}
\leanline{0}{def\ klmcStep\ (force\ :\ E\ →\ E)\ (γ\ h\ :\ ℝ)}
\leanline{4}{(state\ noise\ :\ PhaseSpace\ E)\ :\ PhaseSpace\ E\ :=}
\leanline{2}{exponentialStep\ force\ γ\ h\ state\ (Real.sqrt\ (2\ *\ γ)\ •\ noise)}
\end{leancode}
\leanlabel{How the Lean proof follows the argument}
\texttt{exponential\allowbreak{}Step} records the integrated deterministic friction and frozen-force responses, together with the supplied innovation pair. \texttt{klmc\allowbreak{}Step} substitutes the coefficient functions into that same expression. Their definitions match the two update equations term by term: the first coordinate is velocity and the second is position. Independence and covariance are properties of the innovation law established in the following section, rather than extra terms in this deterministic map.
\end{leanbox}

\section{\texorpdfstring{The joint law of velocity and position innovations}{The joint law of velocity and position innovations}}


The two random increments come from the same Brownian motion:


\[\xi=\int_0^h\psi_0(h-s)dW_s,
\qquad \xi'=\int_0^h\psi_1(h-s)dW_s.\]


For each spatial coordinate their covariance is the two-by-two Gram matrix


\[C=\int_0^h
\begin{pmatrix}\psi_0(t)^2&\psi_0(t)\psi_1(t)\\
\psi_0(t)\psi_1(t)&\psi_1(t)^2\end{pmatrix}dt.\]


The off-diagonal entry is positive: both responses integrate the same Brownian increment with positive kernels. Replacing these increments with
two independent Gaussians changes the algorithm. Different spatial coordinates
and different time intervals use independent Gaussian pairs.




For any real \(a,b\),


\[(a,b)C(a,b)^T=\int_0^h(a\psi_0(t)+b\psi_1(t))^2dt\ge0.\]


This proves positive semidefiniteness. When \(h>0\), the two coefficient
functions are linearly independent, which gives strict positivity. One may
therefore sample a pair using a Cholesky factor \(C=HH^T\) and two independent
standard normals. A complete law proof identifies this construction with the
Brownian integrals and establishes independence from the current state.



For an explicit scalar Gaussian construction, write \(C_{00}=a\), \(C_{10}=b\), and \(C_{11}=d\). Positive definiteness gives \(a>0\) and \(d-b^2/a>0\). With independent standard normals \(G,H\), the pair \((\sqrt aG,(b/\sqrt a)G+\sqrt{d-b^2/a}\,H)\) has exactly covariance \(C\). Its nonzero cross term retains the positive cross covariance.

\begin{leanbox}{Lean: \texttt{klmc\_\allowbreak{}jointLaw}, \texttt{canonicalKLMCInnovation\_\allowbreak{}independent}, \texttt{kernelNoisePath\_\allowbreak{}independent} and 1 more}
\leansource{\texttt{klmc\_\allowbreak{}jointLaw} \quad KineticLangevin/\allowbreak{}BrownianConvolution.lean, lines 415–419}
\begin{leancode}
\leanline{0}{theorem\ klmc\_jointLaw\ (γ\ :\ ℝ)\ (t\ :\ ℝ≥0)\ :}
\leanline{4}{HasLaw\ (fun\ ω\ =>\ WithLp.toLp\ 2\ (fun\ i\ :\ Fin\ 2\ =>}
\leanline{6}{convolutionLp\ (klmcKernels\ γ\ i)\ (continuous\_klmcKernels\ γ\ i)\ t\ ω))}
\leanline{6}{(klmcGaussian\ γ\ t)\ probability}
\end{leancode}
\leansource{\texttt{canonicalKLMCInnovation\_\allowbreak{}independent} \quad KineticLangevin/\allowbreak{}CanonicalVectorSegments.lean, lines 161–164}
\begin{leancode}
\leanline{0}{theorem\ canonicalKLMCInnovation\_independent\ (d\ :\ ℕ)\ (γ\ :\ ℝ)\ (h\ :\ ℝ≥0)\ :}
\leanline{4}{iIndepFun\ (canonicalKLMCInnovation\ d\ γ\ h)\ (ULA.CanonicalBrownian.vectorProbability\ (Fin\ d))}
\end{leancode}
\leansource{\texttt{kernelNoisePath\_\allowbreak{}independent} \quad KineticLangevin/\allowbreak{}Gaussian.lean, lines 126–130}
\begin{leancode}
\leanline{0}{theorem\ kernelNoisePath\_independent\ (d\ :\ ℕ)\ (kernels\ :\ ι\ →\ ℝ\ →\ ℝ)\ (h\ :\ ℝ)\ :}
\leanline{4}{iIndepFun\ (fun\ p\ :\ ℕ\ ×\ Fin\ d\ =>}
\leanline{6}{fun\ ω\ :\ ℕ\ ×\ Fin\ d\ →\ EuclideanSpace\ ℝ\ ι\ =>\ ω\ p)}
\leanline{6}{(kernelNoisePath\ d\ kernels\ h)}
\end{leancode}
\leansource{\texttt{klmc\_\allowbreak{}cholesky\_\allowbreak{}independent\_\allowbreak{}normals} \quad KineticLangevin/\allowbreak{}GaussianFactor.lean, lines 93–99}
\begin{leancode}
\leanline{0}{theorem\ klmc\_cholesky\_independent\_normals\ (γ\ h\ :\ ℝ)\ (hh\ :\ 0<h)\ :}
\leanline{4}{HasLaw\ (fun\ z\ :\ Fin\ 2→ℝ\ =>\ toEuclideanCLM\ ({\leanfallback 𝕜}:=ℝ)\ (choleskyTwo\ (klmcCovariance\ γ\ h))\ (WithLp.toLp\ 2\ z))}
\leanline{6}{(klmcGaussian\ γ\ h)\ (Measure.pi\ (fun\ \_\ :\ Fin\ 2\ =>\ gaussianReal\ 0\ 1))}
\end{leancode}
\leanlabel{How the Lean proof follows the argument}
\texttt{klmc\_\allowbreak{}joint\allowbreak{}Law} identifies the joint Gaussian law of the two Brownian convolutions, including their off-diagonal covariance. \texttt{canonical\allowbreak{}KLMCInnovation\_\allowbreak{}independent} and \texttt{kernel\allowbreak{}Noise\allowbreak{}Path\_\allowbreak{}independent} establish independence across the specified source blocks. \texttt{klmc\_\allowbreak{}cholesky\_\allowbreak{}independent\_\allowbreak{}normals} then gives the same law as an explicit linear image of two independent standard normals. Equality of the complete joint law is what permits replacement of the Brownian-derived innovations by the implemented correlated Gaussian vectors.
\end{leanbox}

\section{\texorpdfstring{The local velocity and position errors}{The local velocity and position errors}}


Compare one frozen-force step with the exact kinetic process started at
the same phase-space state. Their Brownian terms cancel. Gradient Lipschitzness
bounds the velocity defect by


\[\|\Delta v_t\|_{L^2}
\le M\int_0^t\|L_s-L_0\|_{L^2}ds
\le\frac{Mt^2}{2}\sup_{0\le r\le h}\|V_r\|_{L^2}.\]


Position is the integral of velocity, so its defect has an additional time
factor:


\[\|\Delta x_h\|_{L^2}\le\frac{Mh^3}{6}
\sup_{0\le r\le h}\|V_r\|_{L^2}.\]


At stationarity \(\|V_r\|_{L^2}=\sqrt p\). For the auxiliary exact path
started from the numerical state, contraction against the stationary exact
path controls its velocity norm by \(\sqrt p+A_k\), where \(A_k\) is the
current error in transformed coordinates. This prevents the initialization
error from reappearing in every step as a fixed additive penalty.




Under the paper's step restriction, combining the two coordinate defects
gives the transformed local estimate
\(\tfrac34Mh^2(\sqrt p+A_k)\). Constructing the auxiliary process, its
moments and the synchronous Gaussian law is part of the analytic proof behind
this estimate.



The orders follow from the two scalar integrals: \(\int_0^hsds=h^2/2\), whereas \(\int_0^h(h-s)sds=h^3/6\). Smoothness is used only to turn the position displacement into force variation. The reference velocity moment, rather than a bound on Brownian path derivatives, controls the remaining stochastic quantity.

\begin{leanbox}{Lean: \texttt{velocityDefect\_\allowbreak{}bound}, \texttt{positionDefect\_\allowbreak{}bound}, \texttt{globalReferenceFlow\_\allowbreak{}stationary}}
\leansource{\texttt{velocityDefect\_\allowbreak{}bound} \quad KineticLangevin/\allowbreak{}LocalError.lean, lines 45–71}
\begin{leancode}
\leanline{0}{theorem\ velocityDefect\_bound\ (g\ :\ E\ →\ E)\ (M\ :\ ℝ≥0)\ (hg\ :\ LipschitzWith\ M\ g)}
\leanline{4}{(x\ v\ :\ ℝ\ →\ E)\ (γ\ T\ C\ :\ ℝ)\ (hγ\ :\ 0\ ≤\ γ)\ (hC\ :\ 0\ ≤\ C)}
\leanline{4}{(hbound\ :\ ∀\ s\ ∈\ Icc\ 0\ T,\ ‖v\ s‖\ ≤\ C)}
\leanline{4}{(hkin\ :\ ∀\ s\ ∈\ Icc\ 0\ T,\ x\ s\ -\ x\ 0\ =\ ∫\ r\ in\ 0..s,\ v\ r)}
\leanline{4}{\{t\ :\ ℝ\}\ (ht\ :\ t\ ∈\ Icc\ 0\ T)\ :}
\leanline{4}{‖velocityDefect\ g\ x\ γ\ t‖\ ≤\ (M\ :\ ℝ)\ *\ C\ *\ t\ \textasciicircum{}\ 2\ /\ 2}
\end{leancode}
\leansource{\texttt{positionDefect\_\allowbreak{}bound} \quad KineticLangevin/\allowbreak{}LocalError.lean, lines 75–95}
\begin{leancode}
\leanline{0}{theorem\ positionDefect\_bound\ (g\ :\ E\ →\ E)\ (M\ :\ ℝ≥0)\ (hg\ :\ LipschitzWith\ M\ g)}
\leanline{4}{(x\ v\ :\ ℝ\ →\ E)\ (γ\ T\ C\ :\ ℝ)\ (hγ\ :\ 0\ ≤\ γ)\ (hC\ :\ 0\ ≤\ C)}
\leanline{4}{(hbound\ :\ ∀\ s\ ∈\ Icc\ 0\ T,\ ‖v\ s‖\ ≤\ C)}
\leanline{4}{(hkin\ :\ ∀\ s\ ∈\ Icc\ 0\ T,\ x\ s\ -\ x\ 0\ =\ ∫\ r\ in\ 0..s,\ v\ r)}
\leanline{4}{\{t\ :\ ℝ\}\ (ht\ :\ t\ ∈\ Icc\ 0\ T)\ :}
\leanline{4}{‖positionDefect\ g\ x\ γ\ t‖\ ≤\ (M\ :\ ℝ)\ *\ C\ *\ t\ \textasciicircum{}\ 3\ /\ 6}
\end{leancode}
\leansource{\texttt{globalReferenceFlow\_\allowbreak{}stationary} \quad KineticLangevin/\allowbreak{}CanonicalKLMC.lean, lines 12–37}
\begin{leancode}
\leanline{0}{theorem\ globalReferenceFlow\_stationary}
\leanline{4}{(f\ :\ H\ d\ →\ ℝ)\ (m\ :\ ℝ)\ (hm\ :\ 0\ <\ m)\ (M\ γ\ :\ ℝ≥0)}
\leanline{4}{(hf\ :\ Differentiable\ ℝ\ f)\ (hsc\ :\ StrongConvexOn\ univ\ m\ f)}
\leanline{4}{(hLip\ :\ LipschitzWith\ M\ (gradient\ f))}
\leanline{4}{(c\ :\ ULA.Coupling\ μ\ (ULA.gibbsMeasure\ f))}
\leanline{4}{(hπ\ :\ MemLp\ (id\ :\ H\ d\ →\ H\ d)\ 2\ (ULA.gibbsMeasure\ f))\ :}
\leanline{4}{GlobalReferenceStationary\ c\ hπ\ f\ M\ γ\ hLip}
\end{leancode}
\leanlabel{How the Lean proof follows the argument}
\texttt{velocity\allowbreak{}Defect\_\allowbreak{}bound} and \texttt{position\allowbreak{}Defect\_\allowbreak{}bound} bound the two frozen-force remainder integrals in the order in which the two update equations are compared. Gradient Lipschitz continuity bounds force differences by position displacements; integration and the stationary velocity second moment supply the powers of the step. \texttt{global\allowbreak{}Reference\allowbreak{}Flow\_\allowbreak{}stationary} proves that the constructed comparator has the Gaussian-times-Gibbs phase law at each time, so the moment used here is derived for that comparator.
\end{leanbox}

\section{\texorpdfstring{Iteration of the first-order coupling inequality}{Iteration of the first-order coupling inequality}}


Let
\(A_k=\| (\Delta V_k+\gamma\Delta L_k,-\Delta V_k)\|_{L^2}\).
Exact kinetic evolution contracts this transformed norm. A triangle inequality
with the one-step defect gives


\[A_{k+1}\le\left(e^{-mh/\gamma}+\frac34Mh^2\right)A_k
+\frac34Mh^2\sqrt p.\]


For \(0<h\le m/(4\gamma M)\), the quadratic correction is small enough
to retain most of the continuous contraction:


\[A_{k+1}\le qA_k+b,\qquad
q=1-\frac{3mh}{4\gamma},\quad b=\frac34Mh^2\sqrt p.\]


Induction gives the finite geometric sum
\(A_k\le q^kA_0+b\sum_{j=0}^{k-1}q^j\). Because \(0\le q<1\),


\[A_k\le q^kA_0+\frac{Mh\gamma\sqrt p}m.\]


Equal initial velocities give \(A_0=\gamma W_2(\nu_0,\pi)\) for an
optimal initial coupling; approximate couplings and a limiting infimum work as
well. The position projection costs \(\sqrt2/\gamma\), yielding the displayed
KLMC theorem. The recurrence itself is elementary. The substantive stochastic
claim is that the chain admits the coupling from which it was derived. Equality of the explicit fresh-noise recursion and its iterated transition laws uses the earlier Markov-law identification (see \cite[Section~3.33.12]{LangevinLean2026ULASC}).



The geometric sum is finite for every iteration count: \(\sum_{j=0}^{n-1}q^j=(1-q^n)/(1-q)\). Bounding it by \((1-q)^{-1}\) separates transient and persistent errors. Because the local forcing is quadratic and the contraction gap is linear in \(h\), the persistent position-error term is linear in the step. The initial-error coefficient remains outside this sum.

\begin{leanbox}{Lean: \texttt{firstOrder\_\allowbreak{}coupling\_\allowbreak{}recurrence\_\allowbreak{}bound}, \texttt{klmc\_\allowbreak{}wasserstein\_\allowbreak{}bound}}
\leansource{\texttt{firstOrder\_\allowbreak{}coupling\_\allowbreak{}recurrence\_\allowbreak{}bound} \quad KineticLangevin/\allowbreak{}Recurrence.lean, lines 69–96}
\begin{leancode}
\leanline{0}{theorem\ firstOrder\_coupling\_recurrence\_bound\ (e\ :\ ℕ\ →\ ℝ)\ (r\ M\ h\ C\ :\ ℝ)}
\leanline{4}{(hr\ :\ 0\ <\ r)\ (hM\ :\ 0\ ≤\ M)\ (hh\ :\ 0\ <\ h)\ (hC\ :\ 0\ ≤\ C)}
\leanline{4}{(hcurvature\ :\ 2\ *\ r\ \textasciicircum{}\ 2\ ≤\ M)\ (hstep\ :\ 4\ *\ M\ *\ h\ ≤\ r)}
\leanline{4}{(hnonneg\ :\ ∀\ n,\ 0\ ≤\ e\ n)}
\leanline{4}{(hlocal\ :\ ∀\ n,\ e\ (n\ +\ 1)\ ≤\ Real.exp\ (-r\ *\ h)\ *\ e\ n\ +}
\leanline{6}{(3\ /\ 4\ :\ ℝ)\ *\ M\ *\ h\ \textasciicircum{}\ 2\ *\ (C\ +\ e\ n))\ (n\ :\ ℕ)\ :}
\leanline{4}{e\ n\ ≤\ (1\ -\ (3\ /\ 4\ :\ ℝ)\ *\ r\ *\ h)\ \textasciicircum{}\ n\ *\ e\ 0\ +\ M\ *\ h\ *\ C\ /\ r}
\end{leancode}
\leansource{\texttt{klmc\_\allowbreak{}wasserstein\_\allowbreak{}bound} \quad KineticLangevin/\allowbreak{}CanonicalKLMC.lean, lines 41–53}
\begin{leancode}
\leanline{0}{theorem\ klmc\_wasserstein\_bound}
\leanline{4}{(f\ :\ H\ d\ →\ ℝ)\ (m\ :\ ℝ)\ (M\ γ\ :\ ℝ≥0)}
\leanline{4}{(hf\ :\ Differentiable\ ℝ\ f)\ (hsc\ :\ StrongConvexOn\ univ\ m\ f)}
\leanline{4}{(hLip\ :\ LipschitzWith\ M\ (gradient\ f))\ (hmM\ :\ m\ ≤\ (M\ :\ ℝ))\ (hm\ :\ 0\ <\ m)}
\leanline{4}{(hγ\ :\ 0\ <\ (γ\ :\ ℝ))\ (hfriction\ :\ m+(M\ :\ ℝ)\ ≤\ (γ\ :\ ℝ)\textasciicircum{}2)}
\leanline{4}{(h\ :\ ℝ≥0)\ (hh\ :\ 0\ <\ (h\ :\ ℝ))\ (hstep\ :\ 4*(γ\ :\ ℝ)*(M\ :\ ℝ)*(h\ :\ ℝ)\ ≤\ m)}
\leanline{4}{(hfinite\ :\ ULA.wasserstein2\ μ\ (ULA.gibbsMeasure\ f)\ ≠\ ⊤)\ (n\ :\ ℕ)\ :}
\leanline{4}{ULA.wasserstein2Real\ (klmcLaw\ f\ γ\ h\ μ\ n)\ (ULA.gibbsMeasure\ f)\ ≤}
\leanline{6}{Real.sqrt\ 2\ *\ (1-3*m*(h\ :\ ℝ)/(4*(γ\ :\ ℝ)))\textasciicircum{}n\ *\ ULA.wasserstein2Real\ μ\ (ULA.gibbsMeasure\ f)\ +}
\leanline{8}{(M\ :\ ℝ)*(h\ :\ ℝ)*Real.sqrt\ (2*(d\ :\ ℝ))/m}
\end{leancode}
\leanlabel{How the Lean proof follows the argument}
\texttt{first\allowbreak{}Order\_\allowbreak{}coupling\_\allowbreak{}recurrence\_\allowbreak{}bound} combines one-step contraction and local error, iterates the scalar affine recurrence and projects its transformed phase norm to position. \texttt{klmc\_\allowbreak{}wasserstein\_\allowbreak{}bound} obtains the Gibbs second moment as \texttt{h\ensuremath{\pi}}, then applies \texttt{klmc\_\allowbreak{}wasserstein\_\allowbreak{}bound\_\allowbreak{}of\_\allowbreak{}global\allowbreak{}Stationarity} with the stationarity function proved for every initial coupling. The helper theorem identifies both output marginals and takes the initial-coupling infimum. The displayed transient and bias are therefore conclusions for the actual sampler law.
\end{leanbox}

\section{\texorpdfstring{The Hessian-corrected exponential update}{The Hessian-corrected exponential update}}


If the Hessian is \(M_2\)-Lipschitz in operator norm, a first-order
Taylor expansion of the gradient leaves a quadratic position remainder:


\[\|\nabla f(y)-\nabla f(x)-\nabla^2f(x)(y-x)\|
\le\frac{M_2}{2}\|y-x\|^2.\]


The paper inserts the leading velocity and Brownian contributions to
\(L_s-L_0\) into this expansion. Define
\(\varphi_{j+1}(h)=\int_0^he^{-\gamma(h-s)}\psi_j(s)ds\). The resulting
KLMC2 update subtracts \(\varphi_2(h)H v\) from velocity and
\(\varphi_3(h)H v\) from position, where \(H=\nabla^2f(x)\). Its random
increments are also corrected by multiplying two additional correlated
Gaussian blocks by \(H\).




The joint covariance is the Gram integral of
\((\psi_0,\psi_1,\varphi_2,\varphi_3)\). This avoids a matrix exponential:
the operations involving the Hessian are matrix–vector products. The second
algorithm is therefore more elaborate than the frozen-force integrator, but
still explicit.




The name KLMC2 here refers to a Hessian method. KIPLMC2 in the particle
paper uses the OBABO splitting and first derivatives; its definition and proof
are separate.



The frozen Hessian acts on the displacement generated by the free Ornstein–Uhlenbeck velocity. That displacement contains both the initial-velocity response and fresh Brownian responses. Applying the Hessian to both parts is necessary to obtain the stated higher-order approximation; retaining only the deterministic Hessian–velocity correction would define a different algorithm.

\begin{leanbox}{Lean: \texttt{klmc2Step}, \texttt{klmc2\_\allowbreak{}jointLaw}}
\leansource{\texttt{klmc2Step} \quad KineticLangevin/\allowbreak{}Integrator.lean, lines 110–117}
\begin{leancode}
\leanline{0}{def\ klmc2Step\ (force\ :\ E\ →\ E)\ (hessian\ :\ E\ →\ E\ →L[ℝ]\ E)\ (γ\ h\ :\ ℝ)}
\leanline{4}{(state\ :\ PhaseSpace\ E)\ (noise\ :\ Fin\ 4\ →\ E)\ :\ PhaseSpace\ E\ :=}
\leanline{2}{(psi\ γ\ 0\ h\ •\ state.1\ -\ psi\ γ\ 1\ h\ •\ force\ state.2\ -}
\leanline{6}{phi\ γ\ 2\ h\ •\ hessian\ state.2\ state.1\ +}
\leanline{6}{Real.sqrt\ (2\ *\ γ)\ •\ (noise\ 0\ -\ hessian\ state.2\ (noise\ 2)),}
\leanline{3}{state.2\ +\ psi\ γ\ 1\ h\ •\ state.1\ -\ psi\ γ\ 2\ h\ •\ force\ state.2\ -}
\leanline{6}{phi\ γ\ 3\ h\ •\ hessian\ state.2\ state.1\ +}
\leanline{6}{Real.sqrt\ (2\ *\ γ)\ •\ (noise\ 1\ -\ hessian\ state.2\ (noise\ 3)))}
\end{leancode}
\leansource{\texttt{klmc2\_\allowbreak{}jointLaw} \quad KineticLangevin/\allowbreak{}BrownianConvolution.lean, lines 423–427}
\begin{leancode}
\leanline{0}{theorem\ klmc2\_jointLaw\ (γ\ :\ ℝ)\ (t\ :\ ℝ≥0)\ :}
\leanline{4}{HasLaw\ (fun\ ω\ =>\ WithLp.toLp\ 2\ (fun\ i\ :\ Fin\ 4\ =>}
\leanline{6}{convolutionLp\ (klmc2Kernels\ γ\ i)\ (continuous\_klmc2Kernels\ γ\ i)\ t\ ω))}
\leanline{6}{(klmc2Gaussian\ γ\ t)\ probability}
\end{leancode}
\leanlabel{How the Lean proof follows the argument}
\texttt{klmc2\allowbreak{}Step} defines the complete Hessian-corrected update. \texttt{klmc2\_\allowbreak{}joint\allowbreak{}Law} identifies all four Brownian response coordinates jointly. This includes the stochastic responses acted on by the frozen Hessian, as well as the initial-velocity response. The final theorem \texttt{klmc2\_\allowbreak{}wasserstein\_\allowbreak{}bound\_\allowbreak{}exact} concerns this complete update, so removing one of its stochastic Hessian contributions would change the algorithm and invalidate the correspondence.
\end{leanbox}

\section{\texorpdfstring{Fourth moments and stochastic Hessian corrections}{Fourth moments and stochastic Hessian corrections}}


The Taylor remainder contains \(\|L_s-L_0\|^2\). Its
\(L^2\) norm therefore uses fourth moments. Under the stationary
kinetic law, each \(V_s\) is standard Gaussian and


\[\mathbb E\|V_s\|^2=p,\qquad
\mathbb E\|V_s\|^4=p^2+2p.\]


The kinematic equation \(L_s-L_0=\int_0^sV_rdr\) gives
\(\|L_s-L_0\|_{L^4}\le s(p^2+2p)^{1/4}\), without independence
between velocities at different times. The proof applies closed-convex
averaging to the existing \(L^2\) time integral. The Taylor forcing
is consequently bounded by \(M_2s^2\sqrt{p^2+2p}/2\).




Its damped integral gives the cubic velocity defect. A further force-memory
term uses the stationary score bound
\(\mathbb E_\pi\|\nabla f\|^2\le Mp\) from the earlier Gibbs score theorem (see \cite[Section~3.27.8]{LangevinLean2026ULASC}). Integrating the velocity error
once more gives the quartic position defect. Proposition 2 is thus for the defined four-kernel step:


\[\|P^{-1}(V_h-\hat V_h,L_h-\hat L_h)\|_{L^2}
\le\frac{h^3}{4}\bigl(M_2\sqrt{p^2+2p}+M\sqrt{Mp}\bigr),
\qquad \gamma h\le\tfrac15.\]


Proposition 3 requires an additional repair: the source's displayed
subtraction omits the Hessian difference multiplying the fresh
\(\varphi_2\) and \(\varphi_3\) noise. The proof keeps both terms.
Write \(D_H=H(L_0)-H(\widetilde L_0)\), and define the free-OU velocity
and complete Hessian responses on the same Brownian interval:


\[U_r=\psi_0(r)V_0+\sqrt{2\gamma}\int_0^r\psi_0(r-s)dB_s,\]
\[C_j(t)=\varphi_j(t)V_0+\sqrt{2\gamma}
\int_0^t\varphi_j(t-s)dB_s,\qquad j=2,3.\]


Kernel identities, proved by Hilbert-space Fubini and the Brownian
isonormal map, give


\[C_2(t)=\int_0^te^{-\gamma(t-s)}\int_0^sU_rdrds,
\qquad C_2(t)+\gamma C_3(t)=\int_0^t\int_0^sU_rdrds.\]


Each \(U_r\) has the standard Gaussian marginal, proved from the stationary reference velocity and its independent future Brownian interval.
The random operator \(D_H\) can depend on the coupled positions and velocity;
the cutoff argument below requires no independence from \(U_r\). The complete
correction therefore satisfies, for \(a\ge5p\),


\[\|P^{-1}(D_HC_2(h),D_HC_3(h))\|_{L^2}
\le\frac{h^2}{\sqrt2}\left[
M_2\sqrt a\,\|L_0-\widetilde L_0\|_{L^2}
+2(M-m)e^{-(a-p)/8}\right].\]


This retains the source constants while accounting for every stochastic
Hessian term. The frozen-linear part contracts by \(1-mh/(2\gamma)\)
under \(\gamma^2\ge m+M\) and \(5\gamma Mh\le m\); adding the complete
correction gives the repaired one-step comparison.



The fourth moment enters because a quadratic Taylor remainder is estimated in second mean: \(\|\|D\|^2\|_{L^2}=(\mathbb E\|D\|^4)^{1/2}\). Expanding the Gaussian fourth norm gives \(3p+p(p-1)=p^2+2p\). This exact coefficient is kept until any later dimension-dependent simplification is explicitly justified.

\begin{leanbox}{Lean: \texttt{stdGaussian\_\allowbreak{}norm\_\allowbreak{}fourth\_\allowbreak{}moment}, \texttt{klmc2Step\_\allowbreak{}stationary\_\allowbreak{}local\_\allowbreak{}bound}, \texttt{klmc2SourceStep\_\allowbreak{}comparison\_\allowbreak{}bound} and 1 more}
\leansource{\texttt{stdGaussian\_\allowbreak{}norm\_\allowbreak{}fourth\_\allowbreak{}moment} \quad KineticLangevin/\allowbreak{}GaussianFourth.lean, lines 53–106}
\begin{leancode}
\leanline{0}{theorem\ stdGaussian\_norm\_fourth\_moment\ (ι\ :\ Type*)\ [Fintype\ ι]\ :}
\leanline{4}{(∫\ z\ :\ EuclideanSpace\ ℝ\ ι,\ ‖z‖\textasciicircum{}4\ ∂stdGaussian\ (EuclideanSpace\ ℝ\ ι))\ =}
\leanline{6}{(Fintype.card\ ι\ :\ ℝ)\textasciicircum{}2+2*(Fintype.card\ ι\ :\ ℝ)}
\end{leancode}
\leansource{\texttt{klmc2Step\_\allowbreak{}stationary\_\allowbreak{}local\_\allowbreak{}bound} \quad KineticLangevin/\allowbreak{}SecondOrderOneStep.lean, lines 24–50}
\begin{leancode}
\leanline{0}{theorem\ klmc2Step\_stationary\_local\_bound}
\leanline{4}{(f\ :\ EuclideanSpace\ ℝ\ ι\ →\ ℝ)}
\leanline{4}{(H\ :\ EuclideanSpace\ ℝ\ ι\ →\ EuclideanSpace\ ℝ\ ι\ →L[ℝ]\ EuclideanSpace\ ℝ\ ι)}
\leanline{4}{(m\ γ\ T\ :\ ℝ)\ (M\ M2\ :\ ℝ≥0)\ (hm\ :\ 0<m)\ (hγ\ :\ 0≤γ)\ (hT\ :\ 0≤T)}
\leanline{4}{(hsmall\ :\ γ*T≤1/5)}
\leanline{4}{(hf\ :\ Differentiable\ ℝ\ f)\ (hsc\ :\ StrongConvexOn\ univ\ m\ f)}
\leanline{4}{(hLip\ :\ LipschitzWith\ M\ (gradient\ f))\ (hH\ :\ LipschitzWith\ M2\ H)}
\leanline{4}{(hderiv\ :\ ∀\ x,HasFDerivAt\ (gradient\ f)\ (H\ x)\ x)}
\leanline{4}{(R\ :\ Lp\ (EuclideanSpace\ ℝ\ ι)\ 2\ (ULA.CanonicalBrownian.vectorProbability\ ι)\ →L[ℝ]}
\leanline{6}{Lp\ (EuclideanSpace\ ℝ\ ι)\ 2\ P)\ (a\ :\ ℝ≥0)}
\leanline{4}{(V\ X\ :\ ℝ\ →\ Lp\ (EuclideanSpace\ ℝ\ ι)\ 2\ P)\ (hV\ :\ Continuous\ V)\ (hX\ :\ Continuous\ X)}
\leanline{4}{(hlaw\ :\ ∀\ t∈Icc\ 0\ T,P.map\ (lpPhaseCoe\ P\ (V\ t,X\ t))=kineticGibbsTarget\ f)}
\leanline{4}{(hmild\ :\ ∀\ t∈Icc\ 0\ T,V\ t=freeOUResponse\ P\ (Real.sqrt\ (2*γ)\ •\ R)\ (V\ 0)\ γ\ a\ 0\ t-}
\leanline{6}{dampedIntegral\ γ\ (fun\ s\ ↦\ ULA.lpDrift\ P\ (gradient\ f)\ hLip\ (X\ s))\ t)}
\leanline{4}{(hkin\ :\ ∀\ t∈Icc\ 0\ T,X\ t-X\ 0=∫\ r\ in\ 0..t,V\ r)\ :}
\leanline{4}{let\ G\ :=\ ULA.lpDrift\ P\ (gradient\ f)\ hLip}
\leanline{4}{let\ A\ :=\ lpHessian\ P\ (gradient\ f)\ H\ M\ M2\ hLip\ hH\ hderiv}
\leanline{4}{transformNorm\ γ\ 0\ (klmc2Step\ G\ A\ γ\ T\ (V\ 0,X\ 0)\ (liftedKLMC2Noise\ P\ R\ γ\ a\ T)-(V\ T,X\ T))≤}
\leanline{6}{T\textasciicircum{}3/4*((M2\ :\ ℝ)*Real.sqrt\ ((Fintype.card\ ι\ :\ ℝ)\textasciicircum{}2+2*(Fintype.card\ ι\ :\ ℝ))+}
\leanline{8}{(M\ :\ ℝ)*Real.sqrt\ ((M\ :\ ℝ)*(Fintype.card\ ι\ :\ ℝ)))}
\end{leancode}
\leansource{\texttt{klmc2SourceStep\_\allowbreak{}comparison\_\allowbreak{}bound} \quad KineticLangevin/\allowbreak{}SecondOrderComparison.lean, lines 135–169}
\begin{leancode}
\leanline{0}{theorem\ klmc2SourceStep\_comparison\_bound}
\leanline{4}{(f\ :\ EuclideanSpace\ ℝ\ ι\ →\ ℝ)}
\leanline{4}{(H\ :\ EuclideanSpace\ ℝ\ ι\ →\ EuclideanSpace\ ℝ\ ι\ →L[ℝ]\ EuclideanSpace\ ℝ\ ι)}
\leanline{4}{(m\ γ\ t\ α\ :\ ℝ)\ (M\ M2\ :\ ℝ≥0)\ (hm\ :\ 0<m)\ (hmM\ :\ m≤(M\ :\ ℝ))\ (hγ\ :\ 0<γ)}
\leanline{4}{(hfriction\ :\ m+(M\ :\ ℝ)≤γ\textasciicircum{}2)\ (ht\ :\ 0≤t)\ (hstep\ :\ 5*γ*(M\ :\ ℝ)*t≤m)}
\leanline{4}{(hα\ :\ 5*(Fintype.card\ ι\ :\ ℝ)≤α)}
\leanline{4}{(hf\ :\ Differentiable\ ℝ\ f)\ (hsc\ :\ StrongConvexOn\ univ\ m\ f)\ (hLip\ :\ LipschitzWith\ M\ (gradient\ f))}
\leanline{4}{(hH\ :\ LipschitzWith\ M2\ H)\ (hderiv\ :\ ∀\ x,HasFDerivAt\ (gradient\ f)\ (H\ x)\ x)}
\leanline{4}{(L\ :\ Lp\ (EuclideanSpace\ ℝ\ ι)\ 2\ (ULA.CanonicalBrownian.vectorProbability\ ι)\ →L[ℝ]}
\leanline{6}{Lp\ (EuclideanSpace\ ℝ\ ι)\ 2\ P)\ (a\ :\ ℝ≥0)}
\leanline{4}{(z\ w\ :\ Lp\ (EuclideanSpace\ ℝ\ ι)\ 2\ P×Lp\ (EuclideanSpace\ ℝ\ ι)\ 2\ P)}
\leanline{4}{(hGaussian\ :\ ∀\ r∈Icc\ 0\ t,HasLaw\ (freeOUResponse\ P\ L\ z.1\ γ\ a\ 0\ r)\ (stdGaussian\ (EuclideanSpace\ ℝ\ ι))\ P)\ :}
\leanline{4}{let\ G\ :=\ ULA.lpDrift\ P\ (gradient\ f)\ hLip}
\leanline{4}{let\ A\ :=\ lpHessian\ P\ (gradient\ f)\ H\ M\ M2\ hLip\ hH\ hderiv}
\leanline{4}{transformNorm\ γ\ 0\ (klmc2SourceStep\ P\ G\ A\ L\ γ\ a\ t\ z-klmc2SourceStep\ P\ G\ A\ L\ γ\ a\ t\ w)≤}
\leanline{6}{(1-m*t/(2*γ)+(M2\ :\ ℝ)*Real.sqrt\ α*t\textasciicircum{}2/γ)*transformNorm\ γ\ 0\ (z-w)+}
\leanline{8}{Real.sqrt\ 2*t\textasciicircum{}2*((M\ :\ ℝ)-m)*Real.exp\ (-(α-(Fintype.card\ ι\ :\ ℝ))/8)}
\end{leancode}
\leansource{\texttt{referenceChain\_\allowbreak{}freeOUResponse\_\allowbreak{}hasLaw} \quad KineticLangevin/\allowbreak{}CanonicalFreeOU.lean, lines 12–26}
\begin{leancode}
\leanline{0}{theorem\ referenceChain\_freeOUResponse\_hasLaw}
\leanline{4}{(f\ :\ H\ ι\ →\ ℝ)\ (m\ :\ ℝ)\ (hm\ :\ 0<m)\ (M\ γ\ :\ ℝ≥0)}
\leanline{4}{(hf\ :\ Differentiable\ ℝ\ f)\ (hsc\ :\ StrongConvexOn\ univ\ m\ f)\ (hLip\ :\ LipschitzWith\ M\ (gradient\ f))}
\leanline{4}{(c\ :\ ULA.Coupling\ μ\ (ULA.gibbsMeasure\ f))}
\leanline{4}{(hπ\ :\ MemLp\ (id\ :\ H\ ι\ →\ H\ ι)\ 2\ (ULA.gibbsMeasure\ f))}
\leanline{4}{(h\ :\ ℝ≥0)\ (n\ :\ ℕ)\ (r\ :\ ℝ)\ (hr\ :\ 0≤r)\ :}
\leanline{4}{HasLaw\ (freeOUResponse\ (probability\ c)\ (inputLift\ c\ γ)}
\leanline{6}{(referenceChain\ c\ hπ\ f\ M\ γ\ hLip\ h\ n).1\ γ\ ((n\ :\ ℝ≥0)*h)\ 0\ r)}
\leanline{6}{(stdGaussian\ (H\ ι))\ (probability\ c)}
\end{leancode}
\leanlabel{How the Lean proof follows the argument}
\texttt{std\allowbreak{}Gaussian\_\allowbreak{}norm\_\allowbreak{}fourth\_\allowbreak{}moment} supplies the exact dimension factor in the Hessian remainder. \texttt{klmc2\allowbreak{}Step\_\allowbreak{}stationary\_\allowbreak{}local\_\allowbreak{}bound} uses it in the one-step approximation of the stationary flow. \texttt{klmc2\allowbreak{}Source\allowbreak{}Step\_\allowbreak{}comparison\_\allowbreak{}bound} handles the comparison of two numerical states with the complete stochastic Hessian responses. \texttt{reference\allowbreak{}Chain\_\allowbreak{}free\allowbreak{}OUResponse\_\allowbreak{}has\allowbreak{}Law} supplies the Gaussian marginal used in that comparison from the proved joint past/\allowbreak{}future source law; it does not assume independence of the current discrepancy and velocity.
\end{leanbox}

\section{\texorpdfstring{Gaussian truncation and the second-order sampling bound}{Gaussian truncation and the second-order sampling bound}}


A Gaussian reference velocity need not be independent of the accumulated
position error. For any standard Gaussian vector and \(a\ge5p\),
the Gaussian exponential moment gives


\[\|V\|^2\le a+(\|V\|^2-a)_+,\qquad
\mathbb E(\|V\|^2-a)_+\le4e^{-(a-p)/4}.\]


The Hessian difference satisfies both
\(\|H(x)-H(y)\|\le M-m\) and
\(\|H(x)-H(y)\|\le M_2\|x-y\|\). On the bounded-velocity part use
the Lipschitz bound; on the tail use the curvature-width bound. Taking the
square root of the integrated estimate gives


\[\|D_HV\|_{L^2}\le M_2\sqrt a\,\|L_0-\widetilde L_0\|_{L^2}
+2(M-m)e^{-(a-p)/8}.\]


Apply this estimate to each free-OU velocity in the complete
correction, then combine the repaired comparison with the cubic local
defect. Choose


\[a=\left(\frac{m}{4M_2h}\right)^2,\qquad
h\le\frac{m}{5\gamma M}\wedge\frac{m}{4M_2\sqrt{5p}}.\]


Assume \(m,M_2>0\), \(\gamma^2\ge m+M\), and initial phase law
\(\mathcal N(0,I_p)\otimes\nu_0\), with finite \(W_2(\nu_0,\pi)\).
The second step restriction ensures \(a\ge5p\), and the first enforces the
frozen-linear stability bound. The position-dependent comparison term is
absorbed into \(q=1-mh/(4\gamma)\). Iterating on the constructed common
Brownian source and taking the infimum over initial couplings proves
for the independently initialized, iid four-kernel Gaussian KLMC2 chain


\[W_2(\nu_n^{\rm KLMC2},\pi)\le
\sqrt2q^n W_2(\nu_0,\pi)
+\frac{\sqrt2h^2}{m}\left(M_2\sqrt{p^2+2p}+M\sqrt{Mp}\right)
+\frac{8(M-m)h}{m}e^{-(a-p)/8}.\]


The output transport distance is finite; the real-valued expression
does not conceal an infinite distance. No stationary-law, local-error,
independence, Gaussian-innovation or optimal-coupling premise is supplied to
this final sampler theorem.




The cutoff gives both \(e^{-(a-p)/8}\le e^{-p/2}\) and the stronger
\(e^{-(a-p)/8}\le e^{-m^2/(160M_2^2h^2)}\). The printed simplification
of the first bias term uses \(p^2+2p\le2p^2\).




\textbf{Extra assumption for that simplification:} \(p\ge2\).




Under this condition the bound implies


\[W_2(\nu_n^{\rm KLMC2},\pi)\le
\sqrt2q^n W_2(\nu_0,\pi)
+\frac{2h^2M_2p}{m}+\frac{h^2M\sqrt{2Mp}}m
+\frac{8Mh}{m}e^{-p/2}.\]


For \(p=1\), the exact fourth-moment expression in the preceding
sampling theorem remains the valid conclusion. The source's silently used
higher-dimensional simplification is not applied.



The cutoff gives two estimates. Its moderate-velocity contribution is proportional to the current displacement and is absorbed into the contraction gap. Its remaining excess is controlled solely by the Gaussian velocity marginal. Thus the argument remains valid when the displacement and velocity share the entire past of the coupled sampler.

\begin{leanbox}{Lean: \texttt{gaussian\_\allowbreak{}norm\_\allowbreak{}square\_\allowbreak{}positivePart\_\allowbreak{}le}, \texttt{gaussian\_\allowbreak{}truncation\_\allowbreak{}Lp}, \texttt{secondOrderCutoff\_\allowbreak{}tail} and 2 more}
\leansource{\texttt{gaussian\_\allowbreak{}norm\_\allowbreak{}square\_\allowbreak{}positivePart\_\allowbreak{}le} \quad KineticLangevin/\allowbreak{}GaussianTail.lean, lines 124–131}
\begin{leancode}
\leanline{0}{theorem\ gaussian\_norm\_square\_positivePart\_le\ \{Ω\ ι\ :\ Type*\}\ [MeasurableSpace\ Ω]\ [Fintype\ ι]}
\leanline{4}{(P\ :\ Measure\ Ω)\ (V\ :\ Ω\ →\ EuclideanSpace\ ℝ\ ι)}
\leanline{4}{(hV\ :\ HasLaw\ V\ (stdGaussian\ (EuclideanSpace\ ℝ\ ι))\ P)}
\leanline{4}{(a\ :\ ℝ)\ (ha\ :\ 5*(Fintype.card\ ι\ :\ ℝ)≤a)\ :}
\leanline{4}{(∫\ ω,max\ (‖V\ ω‖\textasciicircum{}2-a)\ 0\ ∂P)≤4*Real.exp\ (-(a-(Fintype.card\ ι\ :\ ℝ))/4)}
\end{leancode}
\leansource{\texttt{gaussian\_\allowbreak{}truncation\_\allowbreak{}Lp} \quad KineticLangevin/\allowbreak{}GaussianTruncation.lean, lines 55–96}
\begin{leancode}
\leanline{0}{theorem\ gaussian\_truncation\_Lp\ \{Ω\ ι\ :\ Type*\}\ [MeasurableSpace\ Ω]\ [Fintype\ ι]}
\leanline{4}{(P\ :\ Measure\ Ω)\ [IsProbabilityMeasure\ P]}
\leanline{4}{(D\ X\ V\ :\ Lp\ (EuclideanSpace\ ℝ\ ι)\ 2\ P)}
\leanline{4}{(hV\ :\ HasLaw\ V\ (stdGaussian\ (EuclideanSpace\ ℝ\ ι))\ P)}
\leanline{4}{(B\ L\ a\ :\ ℝ)\ (hB\ :\ 0≤B)\ (hL\ :\ 0≤L)\ (ha\ :\ 5*(Fintype.card\ ι\ :\ ℝ)≤a)}
\leanline{4}{(hbound\ :\ ∀ᵐ\ ω\ ∂P,‖D\ ω‖≤min\ B\ (L*‖X\ ω‖)*‖V\ ω‖)\ :}
\leanline{4}{‖D‖≤L*Real.sqrt\ a*‖X‖+2*B*Real.exp\ (-(a-(Fintype.card\ ι\ :\ ℝ))/8)}
\end{leancode}
\leansource{\texttt{secondOrderCutoff\_\allowbreak{}tail} \quad KineticLangevin/\allowbreak{}SecondOrderRecurrence.lean, lines 27–38}
\begin{leancode}
\leanline{0}{theorem\ secondOrderCutoff\_tail\ (m\ M2\ h\ p\ :\ ℝ)\ (hm\ :\ 0<m)\ (hM2\ :\ 0<M2)}
\leanline{4}{(hh\ :\ 0<h)\ (hp\ :\ 0≤p)\ (hdim\ :\ 4*Real.sqrt\ (5*p)*M2*h≤m)\ :}
\leanline{4}{Real.exp\ (-(secondOrderCutoff\ m\ M2\ h-p)/8)≤Real.exp\ (-p/2)\ ∧}
\leanline{4}{Real.exp\ (-(secondOrderCutoff\ m\ M2\ h-p)/8)≤Real.exp\ (-m\textasciicircum{}2/(160*M2\textasciicircum{}2*h\textasciicircum{}2))}
\end{leancode}
\leansource{\texttt{klmc2\_\allowbreak{}wasserstein\_\allowbreak{}bound\_\allowbreak{}exact} \quad KineticLangevin/\allowbreak{}CanonicalKLMC2.lean, lines 13–62}
\begin{leancode}
\leanline{0}{theorem\ klmc2\_wasserstein\_bound\_exact\ \{μ\ :\ Measure\ (H\ ι)\}\ [IsProbabilityMeasure\ μ]}
\leanline{4}{(f\ :\ H\ ι→ℝ)\ (Hess\ :\ H\ ι→H\ ι→L[ℝ]\ H\ ι)\ (m\ :\ ℝ)\ (M\ γ\ M2\ :\ ℝ≥0)}
\leanline{4}{(hf\ :\ Differentiable\ ℝ\ f)\ (hsc\ :\ StrongConvexOn\ univ\ m\ f)}
\leanline{4}{(hLip\ :\ LipschitzWith\ M\ (gradient\ f))\ (hHLip\ :\ LipschitzWith\ M2\ Hess)}
\leanline{4}{(hderiv\ :\ ∀\ x,HasFDerivAt\ (gradient\ f)\ (Hess\ x)\ x)}
\leanline{4}{(hm\ :\ 0<m)\ (hmM\ :\ m≤(M:ℝ))\ (hγ\ :\ 0<(γ:ℝ))\ (hM2\ :\ 0<(M2:ℝ))}
\leanline{4}{(hfriction\ :\ m+(M:ℝ)≤(γ:ℝ)\textasciicircum{}2)\ (h\ :\ ℝ≥0)\ (hh\ :\ 0<(h:ℝ))}
\leanline{4}{(hstep\ :\ 5*(γ:ℝ)*(M:ℝ)*(h:ℝ)≤m)}
\leanline{4}{(hdim\ :\ 4*Real.sqrt\ (5*(Fintype.card\ ι:ℝ))*(M2:ℝ)*(h:ℝ)≤m)}
\leanline{4}{(hfinite\ :\ ULA.wasserstein2\ μ\ (ULA.gibbsMeasure\ f)≠⊤)\ (n\ :\ ℕ)\ :}
\leanline{4}{ULA.wasserstein2Real\ (law\ ι\ (gradient\ f)\ Hess\ γ\ h\ μ\ n)\ (ULA.gibbsMeasure\ f)≤}
\leanline{6}{Real.sqrt\ 2*(1-m*(h:ℝ)/(4*(γ:ℝ)))\textasciicircum{}n*ULA.wasserstein2Real\ μ\ (ULA.gibbsMeasure\ f)+}
\leanline{6}{Real.sqrt\ 2*(h:ℝ)\textasciicircum{}2*((M2:ℝ)*Real.sqrt\ ((Fintype.card\ ι:ℝ)\textasciicircum{}2+2*(Fintype.card\ ι:ℝ))+}
\leanline{8}{(M:ℝ)*Real.sqrt\ ((M:ℝ)*(Fintype.card\ ι:ℝ)))/m+}
\leanline{6}{8*((M:ℝ)-m)*(h:ℝ)/m*Real.exp\ (-(secondOrderCutoff\ m\ M2\ h-(Fintype.card\ ι:ℝ))/8)}
\end{leancode}
\leansource{\texttt{klmc2\_\allowbreak{}wasserstein\_\allowbreak{}bound} \quad KineticLangevin/\allowbreak{}CanonicalKLMC2.lean, lines 88–122}
\begin{leancode}
\leanline{0}{theorem\ klmc2\_wasserstein\_bound\ \{μ\ :\ Measure\ (H\ ι)\}\ [IsProbabilityMeasure\ μ]}
\leanline{4}{(hp\ :\ 2≤Fintype.card\ ι)}
\leanline{4}{(f\ :\ H\ ι→ℝ)\ (Hess\ :\ H\ ι→H\ ι→L[ℝ]\ H\ ι)\ (m\ :\ ℝ)\ (M\ γ\ M2\ :\ ℝ≥0)}
\leanline{4}{(hf\ :\ Differentiable\ ℝ\ f)\ (hsc\ :\ StrongConvexOn\ univ\ m\ f)}
\leanline{4}{(hLip\ :\ LipschitzWith\ M\ (gradient\ f))\ (hHLip\ :\ LipschitzWith\ M2\ Hess)}
\leanline{4}{(hderiv\ :\ ∀\ x,HasFDerivAt\ (gradient\ f)\ (Hess\ x)\ x)}
\leanline{4}{(hm\ :\ 0<m)\ (hmM\ :\ m≤(M:ℝ))\ (hγ\ :\ 0<(γ:ℝ))\ (hM2\ :\ 0<(M2:ℝ))}
\leanline{4}{(hfriction\ :\ m+(M:ℝ)≤(γ:ℝ)\textasciicircum{}2)\ (h\ :\ ℝ≥0)\ (hh\ :\ 0<(h:ℝ))}
\leanline{4}{(hstep\ :\ 5*(γ:ℝ)*(M:ℝ)*(h:ℝ)≤m)}
\leanline{4}{(hdim\ :\ 4*Real.sqrt\ (5*(Fintype.card\ ι:ℝ))*(M2:ℝ)*(h:ℝ)≤m)}
\leanline{4}{(hfinite\ :\ ULA.wasserstein2\ μ\ (ULA.gibbsMeasure\ f)≠⊤)\ (n\ :\ ℕ)\ :}
\leanline{4}{ULA.wasserstein2Real\ (law\ ι\ (gradient\ f)\ Hess\ γ\ h\ μ\ n)\ (ULA.gibbsMeasure\ f)≤}
\leanline{6}{Real.sqrt\ 2*(1-m*(h:ℝ)/(4*(γ:ℝ)))\textasciicircum{}n*ULA.wasserstein2Real\ μ\ (ULA.gibbsMeasure\ f)+}
\leanline{6}{2*(h:ℝ)\textasciicircum{}2*(M2:ℝ)*(Fintype.card\ ι:ℝ)/m+}
\leanline{6}{(h:ℝ)\textasciicircum{}2*(M:ℝ)*Real.sqrt\ (2*(M:ℝ)*(Fintype.card\ ι:ℝ))/m+}
\leanline{6}{(8*(M:ℝ)/m)*(h:ℝ)*Real.exp\ (-(Fintype.card\ ι:ℝ)/2)}
\end{leancode}
\leanlabel{How the Lean proof follows the argument}
\texttt{gaussian\_\allowbreak{}norm\_\allowbreak{}square\_\allowbreak{}positive\allowbreak{}Part\_\allowbreak{}le} gives the Gaussian excess-moment bound, and \texttt{gaussian\_\allowbreak{}truncation\_\allowbreak{}Lp} converts it into the norm estimate needed by the cutoff comparison. \texttt{second\allowbreak{}Order\allowbreak{}Cutoff\_\allowbreak{}tail} evaluates the chosen cutoff and gives both the dimension and inverse-step-square tail bounds. \texttt{klmc2\_\allowbreak{}wasserstein\_\allowbreak{}bound\_\allowbreak{}exact} retains the exact fourth moment. \texttt{klmc2\_\allowbreak{}wasserstein\_\allowbreak{}bound} subsequently enlarges it to the printed coefficient under its explicit dimension hypothesis; the exact theorem remains available in dimension one.
\end{leanbox}

\section{\texorpdfstring{Step restrictions and iteration counts for prescribed accuracy}{Step restrictions and iteration counts for prescribed accuracy}}


Suppose an error bound has the form \(q^kD+Bh\). Choose
\(h\) small enough that \(Bh\) is bounded by one part of the error tolerance
\(\varepsilon\), and choose \(k\) large enough for the decaying initial
term to satisfy the remaining tolerance. The admissible \(h\) is the minimum of this accuracy
choice and the stability restriction.




For the first integrator, the bias scales as \(h\sqrt p\), giving
\(h\) of order \(\varepsilon/\sqrt p\), and a step count of order
\(\sqrt p/\varepsilon\) up to logarithmic and curvature factors.
For the Hessian method the leading bias scales as \(h^2p\), giving
\(h\) of order \(\sqrt{\varepsilon/p}\), and a step count of order
\(\sqrt{p/\varepsilon}\). These are asymptotic consequences of algorithmic estimates, with fixed geometry constants.




The paper's comparisons keep track of condition numbers as well. A kinetic
method gains from smooth position paths, but poor conditioning can make its
allowed step size small. The resulting iteration bound therefore depends on dimension, error tolerance, and curvature. A precise
complexity theorem also specifies integer rounding and the treatment of
already-small initialization error.




The source also gives a finer condition-number comparison after explicitly
omitting the Gaussian tail. We can audit that calculation using exact finite expressions. For target error \(\varepsilon/\sqrt m\), allocate
\(\varepsilon/(2\sqrt m)\) to the printed quadratic bias and set


\[h_0=\min\left\{\frac m{5\gamma M},\;
\frac m{4\sqrt{5p}M_2},\;
\sqrt{\frac{\varepsilon\sqrt m}{4M_2p+2M\sqrt{2Mp}}}\right\}.\]


Every denominator is positive. This step meets both source stability
restrictions and the stated quadratic-bias allocation. Its reciprocal is
exactly


\[Q=\max\left\{\frac{5\gamma M}m,\;
\frac{4\sqrt{5p}M_2}m,\;
\sqrt{\frac{4M_2p+2M\sqrt{2Mp}}{\varepsilon\sqrt m}}\right\}.\]


There is no hidden dimension or curvature factor in this identity. For
initial distance \(D\), the rounded logarithmic iteration prescription satisfies


\[n\le1+\frac{4\gamma}{m}Q\,
\max\left\{0,\log\frac{24D\sqrt m}{\varepsilon}\right\}.\]


This calculation concerns the source comparison with the tail omitted.
The sampler accuracy theorem retains the tail and chooses its step
using the complete quadratic envelope proved earlier.




To understand the refinement, write
\(\kappa=M/m\), \(\eta=M_2/m^{3/2}\), and


\[\rho=\kappa_2=\eta^{2/3}+\kappa p^{-1/3}
=\frac{M_2^{2/3}+Mp^{-1/3}}m.\]


With \(\gamma=\sqrt{m+M}\), the source tries to simplify the count
coefficient to \(\sqrt p\,\rho^2+\sqrt{p/\varepsilon}\,\rho^{5/4}\).
That simplification drops a factor. A valid refinement is


\[\sqrt{\frac\kappa\rho}\left[
\sqrt p\,\rho^2+\sqrt{p/\varepsilon}\,\rho^{5/4}\right]\]


This comparison holds up to a universal multiplicative constant and the displayed logarithm.
A convenient uniform upper bound uses
\(\widehat\kappa_2=\max\{\kappa,\rho\}\):


\[\mathcal C\le64\left[\sqrt p\,\widehat\kappa_2^2+
\sqrt{p/\varepsilon}\,\widehat\kappa_2^{5/4}\right],\]


Here \(\mathcal C=(4\gamma/m)Q\) is the exact count coefficient
above. The natural count is therefore less than one plus this bound times
the visible nonnegative logarithm.




\textbf{Extra assumption for the printed shape:} \(\rho\ge\kappa\), equivalently \(M(1-p^{-1/3})\le M_2^{2/3}\).




Then the maximum is simply \(\rho\), so the printed shape follows with the same explicit constant.




The distinction matters because the modified number discounts gradient
curvature by \(p^{-1/3}\), while the stability restriction still
contains \(M/m\). Take a natural \(t\ge2\), dimension \(p=t^6\),
and the valid constants \(m=1\), \(M=t^2\), \(M_2=1\),
\(\varepsilon=1\). A diagonal quadratic potential can have smallest
curvature one and largest curvature \(t^2\); its constant Hessian satisfies
the stated Hessian-Lipschitz bound. Here \(\kappa=t^2\),
\(\eta=1\), and \(\rho=2\) exactly. The normalized iteration
coefficient is at least \(20t^4\), whereas the source expression is at most
\(8t^3\). No universal constant multiplying that expression can cover all
natural \(t\). This diagnoses the source's simplification of the sufficient step and iteration bound;
it asserts no lower bound on the best possible behavior of the algorithm.



All inverse-step branches must be retained when deriving a uniform work bound. Inverting a minimum of positive admissible steps produces a maximum of the corresponding positive reciprocal costs. The dominant branch may change with curvature, dimension, Hessian regularity, and requested precision. The integer ceiling then adds less than one iteration to the real logarithmic threshold.

\begin{leanbox}{Lean: \texttt{klmc\_\allowbreak{}accuracy\_\allowbreak{}explicit}, \texttt{firstOrderIterationCount\_\allowbreak{}spec}, \texttt{klmc2\_\allowbreak{}accuracy\_\allowbreak{}explicit} and 6 more}
\leansource{\texttt{klmc\_\allowbreak{}accuracy\_\allowbreak{}explicit} \quad KineticLangevin/\allowbreak{}Accuracy.lean, lines 63–77}
\begin{leancode}
\leanline{0}{theorem\ klmc\_accuracy\_explicit\ \{d\ :\ ℕ\}\ (hd\ :\ 0<d)\ \{μ\ :\ Measure\ (CanonicalCoupling.H\ d)\}}
\leanline{4}{[IsProbabilityMeasure\ μ]\ (f\ :\ CanonicalCoupling.H\ d\ →\ ℝ)\ (m\ :\ ℝ)\ (M\ γ\ :\ ℝ≥0)}
\leanline{4}{(hf\ :\ Differentiable\ ℝ\ f)\ (hsc\ :\ StrongConvexOn\ univ\ m\ f)\ (hLip\ :\ LipschitzWith\ M\ (gradient\ f))}
\leanline{4}{(hmM\ :\ m≤(M:ℝ))\ (hm\ :\ 0<m)\ (hγ\ :\ 0<(γ:ℝ))\ (hfriction\ :\ m+(M:ℝ)≤(γ:ℝ)\textasciicircum{}2)}
\leanline{4}{(hfinite\ :\ ULA.wasserstein2\ μ\ (ULA.gibbsMeasure\ f)≠⊤)\ (ε\ :\ ℝ)\ (hε\ :\ 0<ε)\ :}
\leanline{4}{let\ hs\ :=\ klmcAccuracyStep\_spec\ m\ M\ γ\ d\ ε\ hm\ (hm.trans\_le\ hmM)\ hγ\ (by\ exact\_mod\_cast\ hd)\ hε}
\leanline{4}{let\ h\ :\ ℝ≥0\ :=\ ⟨klmcAccuracyStep\ m\ M\ γ\ d\ ε,hs.1.le⟩}
\leanline{4}{let\ n\ :=\ firstOrderIterationCount\ (3*m*(h:ℝ)/(4*(γ:ℝ)))\ (ULA.wasserstein2Real\ μ\ (ULA.gibbsMeasure\ f))\ ε}
\leanline{4}{ULA.wasserstein2Real\ (CanonicalCoupling.klmcLaw\ f\ γ\ h\ μ\ n)\ (ULA.gibbsMeasure\ f)≤ε}
\end{leancode}
\leansource{\texttt{firstOrderIterationCount\_\allowbreak{}spec} \quad KineticLangevin/\allowbreak{}Complexity.lean, lines 33–41}
\begin{leancode}
\leanline{0}{theorem\ firstOrderIterationCount\_spec\ (α\ D\ ε\ :\ ℝ)\ (hα\ :\ 0\ <\ α)\ (n\ :\ ℕ)}
\leanline{4}{(hn\ :\ firstOrderIterationCount\ α\ D\ ε\ ≤\ n)\ :}
\leanline{4}{Real.log\ (24*D/ε)\ ≤\ α*(n\ :\ ℝ)}
\end{leancode}
\leansource{\texttt{klmc2\_\allowbreak{}accuracy\_\allowbreak{}explicit} \quad KineticLangevin/\allowbreak{}SecondOrderAccuracy.lean, lines 35–55}
\begin{leancode}
\leanline{0}{theorem\ klmc2\_accuracy\_explicit\ \{μ\ :\ Measure\ (H\ ι)\}\ [IsProbabilityMeasure\ μ]}
\leanline{4}{(hp\ :\ 0<Fintype.card\ ι)}
\leanline{4}{(f\ :\ H\ ι→ℝ)\ (Hess\ :\ H\ ι→H\ ι→L[ℝ]\ H\ ι)\ (m\ :\ ℝ)\ (M\ γ\ M2\ :\ ℝ≥0)}
\leanline{4}{(hf\ :\ Differentiable\ ℝ\ f)\ (hsc\ :\ StrongConvexOn\ univ\ m\ f)}
\leanline{4}{(hLip\ :\ LipschitzWith\ M\ (gradient\ f))\ (hHLip\ :\ LipschitzWith\ M2\ Hess)}
\leanline{4}{(hderiv\ :\ ∀\ x,HasFDerivAt\ (gradient\ f)\ (Hess\ x)\ x)}
\leanline{4}{(hm\ :\ 0<m)\ (hmM\ :\ m≤(M:ℝ))\ (hγ\ :\ 0<(γ:ℝ))\ (hM2\ :\ 0<(M2:ℝ))}
\leanline{4}{(hfriction\ :\ m+(M:ℝ)≤(γ:ℝ)\textasciicircum{}2)}
\leanline{4}{(hfinite\ :\ ULA.wasserstein2\ μ\ (ULA.gibbsMeasure\ f)≠⊤)\ (ε\ :\ ℝ)\ (hε\ :\ 0<ε)\ :}
\leanline{4}{let\ hs\ :=\ secondOrderAccuracyStep\_spec\ m\ M\ γ\ M2\ (Fintype.card\ ι)\ ε\ hm\ (hm.trans\_le\ hmM)}
\leanline{6}{hγ\ hM2\ (by\ exact\_mod\_cast\ hp)\ hε}
\leanline{4}{let\ h\ :\ ℝ≥0\ :=\ ⟨secondOrderAccuracyStep\ m\ M\ γ\ M2\ (Fintype.card\ ι)\ ε,hs.1.le⟩}
\leanline{4}{let\ n\ :=\ firstOrderIterationCount\ (m*(h:ℝ)/(4*(γ:ℝ)))\ (ULA.wasserstein2Real\ μ\ (ULA.gibbsMeasure\ f))\ ε}
\leanline{4}{ULA.wasserstein2Real\ (law\ ι\ (gradient\ f)\ Hess\ γ\ h\ μ\ n)\ (ULA.gibbsMeasure\ f)≤ε}
\end{leancode}
\leansource{\texttt{secondOrderAccuracyStep\_\allowbreak{}spec} \quad KineticLangevin/\allowbreak{}SecondOrderBudget.lean, lines 54–79}
\begin{leancode}
\leanline{0}{theorem\ secondOrderAccuracyStep\_spec\ (m\ M\ γ\ M2\ p\ ε\ :\ ℝ)\ (hm\ :\ 0<m)\ (hM\ :\ 0<M)}
\leanline{4}{(hγ\ :\ 0<γ)\ (hM2\ :\ 0<M2)\ (hp\ :\ 0<p)\ (hε\ :\ 0<ε)\ :}
\leanline{4}{0<secondOrderAccuracyStep\ m\ M\ γ\ M2\ p\ ε\ ∧\ secondOrderAccuracyStep\ m\ M\ γ\ M2\ p\ ε≤1\ ∧}
\leanline{4}{5*γ*M*secondOrderAccuracyStep\ m\ M\ γ\ M2\ p\ ε≤m\ ∧}
\leanline{4}{4*Real.sqrt\ (5*p)*M2*secondOrderAccuracyStep\ m\ M\ γ\ M2\ p\ ε≤m\ ∧}
\leanline{4}{secondOrderBiasCoefficient\ m\ M\ M2\ p*(secondOrderAccuracyStep\ m\ M\ γ\ M2\ p\ ε)\textasciicircum{}2≤(47/50:ℝ)*ε}
\end{leancode}
\leansource{\texttt{secondOrderTailFreeStep\_\allowbreak{}spec} \quad KineticLangevin/\allowbreak{}SecondOrderComplexity.lean, lines 21–57}
\begin{leancode}
\leanline{0}{theorem\ secondOrderTailFreeStep\_spec\ (m\ M\ γ\ M2\ p\ ε\ :\ ℝ)}
\leanline{4}{(hm\ :\ 0<m)\ (hM\ :\ 0<M)\ (hγ\ :\ 0<γ)\ (hM2\ :\ 0<M2)\ (hp\ :\ 0<p)\ (hε\ :\ 0<ε)\ :}
\leanline{4}{0<secondOrderTailFreeStep\ m\ M\ γ\ M2\ p\ ε\ ∧}
\leanline{4}{5*γ*M*secondOrderTailFreeStep\ m\ M\ γ\ M2\ p\ ε≤m\ ∧}
\leanline{4}{4*Real.sqrt\ (5*p)*M2*secondOrderTailFreeStep\ m\ M\ γ\ M2\ p\ ε≤m\ ∧}
\leanline{4}{2*(secondOrderTailFreeStep\ m\ M\ γ\ M2\ p\ ε)\textasciicircum{}2*M2*p/m+}
\leanline{6}{(secondOrderTailFreeStep\ m\ M\ γ\ M2\ p\ ε)\textasciicircum{}2*M*Real.sqrt\ (2*M*p)/m≤ε/(2*Real.sqrt\ m)}
\end{leancode}
\leansource{\texttt{secondOrderTailFreeStep\_\allowbreak{}inverse} \quad KineticLangevin/\allowbreak{}SecondOrderComplexity.lean, lines 71–77}
\begin{leancode}
\leanline{0}{theorem\ secondOrderTailFreeStep\_inverse\ (m\ M\ γ\ M2\ p\ ε\ :\ ℝ)}
\leanline{4}{(hm\ :\ 0<m)\ (hM\ :\ 0<M)\ (hγ\ :\ 0<γ)\ (hM2\ :\ 0<M2)\ (hp\ :\ 0<p)\ (hε\ :\ 0<ε)\ :}
\leanline{4}{(secondOrderTailFreeStep\ m\ M\ γ\ M2\ p\ ε)⁻¹=secondOrderTailFreeInverse\ m\ M\ γ\ M2\ p\ ε}
\end{leancode}
\leansource{\texttt{secondOrderTailFree\_\allowbreak{}iterationCount\_\allowbreak{}conditionNumber} \quad KineticLangevin/\allowbreak{}SecondOrderComplexity.lean, lines 177–191}
\begin{leancode}
\leanline{0}{theorem\ secondOrderTailFree\_iterationCount\_conditionNumber\ (m\ M\ M2\ p\ ε\ D\ :\ ℝ)}
\leanline{4}{(hm\ :\ 0<m)\ (hmM\ :\ m≤M)\ (hM2\ :\ 0<M2)\ (hp\ :\ 1≤p)\ (hε\ :\ 0<ε)\ :}
\leanline{4}{(firstOrderIterationCount}
\leanline{6}{(m*secondOrderTailFreeStep\ m\ M\ (Real.sqrt\ (m+M))\ M2\ p\ ε/(4*Real.sqrt\ (m+M)))\ D}
\leanline{6}{(ε/Real.sqrt\ m):ℝ)<}
\leanline{6}{1+64*secondOrderRefinement\ (max\ (kineticConditionNumber\ m\ M)}
\leanline{8}{(sourceHessianConditionNumber\ m\ M\ M2\ p))\ p\ ε*}
\leanline{10}{max\ 0\ (Real.log\ (24*D*Real.sqrt\ m/ε))}
\end{leancode}
\leansource{\texttt{secondOrderNormalizedWork\_\allowbreak{}physical\_\allowbreak{}regime} \quad KineticLangevin/\allowbreak{}ConditionNumber.lean, lines 287–293}
\begin{leancode}
\leanline{0}{theorem\ secondOrderNormalizedWork\_physical\_regime\ (m\ M\ M2\ p\ ε\ :\ ℝ)}
\leanline{4}{(hm\ :\ 0<m)\ (hmM\ :\ m≤M)\ (hM2\ :\ 0≤M2)\ (hp\ :\ 1≤p)\ (hε\ :\ 0<ε)}
\leanline{4}{(hregime\ :\ M*(1-p\ \textasciicircum{}\ (-1/3\ :\ ℝ))≤M2\ \textasciicircum{}\ (2/3\ :\ ℝ))\ :}
\leanline{4}{secondOrderNormalizedWork\ (M/m)\ (M2/(m*Real.sqrt\ m))\ p\ ε≤}
\leanline{6}{64*secondOrderRefinement\ (sourceHessianConditionNumber\ m\ M\ M2\ p)\ p\ ε}
\end{leancode}
\leansource{\texttt{secondOrderSourceRefinement\_\allowbreak{}no\_\allowbreak{}uniform\_\allowbreak{}constant} \quad KineticLangevin/\allowbreak{}ConditionNumber.lean, lines 297–305}
\begin{leancode}
\leanline{0}{theorem\ secondOrderSourceRefinement\_no\_uniform\_constant\ (C\ :\ ℝ)\ (hC\ :\ 0<C)\ :}
\leanline{4}{∃\ n\ :\ ℕ,\ 1≤n\ ∧}
\leanline{6}{C*secondOrderRefinement}
\leanline{8}{(sourceHessianConditionNumber\ 1\ ((n:ℝ)\textasciicircum{}2)\ 1\ ((n:ℝ)\textasciicircum{}6))\ ((n:ℝ)\textasciicircum{}6)\ 1\ <}
\leanline{6}{secondOrderNormalizedWork\ ((n:ℝ)\textasciicircum{}2)\ 1\ ((n:ℝ)\textasciicircum{}6)\ 1}
\end{leancode}
\leanlabel{How the Lean proof follows the argument}
\texttt{klmc\_\allowbreak{}accuracy\_\allowbreak{}explicit} applies the proved sampler estimate to a positive minimum step and a rounded logarithmic count. \texttt{first\allowbreak{}Order\allowbreak{}Iteration\allowbreak{}Count\_\allowbreak{}spec} verifies the count's transient budget. For the second-order method, \texttt{second\allowbreak{}Order\allowbreak{}Accuracy\allowbreak{}Step\_\allowbreak{}spec} verifies the algebraic and tail allocations at the chosen step, and \texttt{klmc2\_\allowbreak{}accuracy\_\allowbreak{}explicit} combines those with its rounded iteration count. Each allocation is checked before the final real-arithmetic addition of the error budgets.
\end{leanbox}

\begin{thebibliography}{ULA-LSI}
\bibitem[ULA-SC]{LangevinLean2026ULASC}
LangevinLean Project. \emph{Unadjusted Langevin under Strong Convexity}. LL-01, 2026. Draft report, content revision d34ead06987c. \url{https://odakyildiz.github.io/LangevinLean/docs/pdf/UnadjustedLangevin/FullProof.pdf}.
\end{thebibliography}

\end{document}
